\chapter{Methods}\label{ch:methods}

The preceding chapters established that no existing tool jointly optimizes map styles and tile data, and introduced the compiler, database, and encoding techniques whose analogues make such joint optimization possible.
This chapter applies those foundations in a concrete Rust optimizer that consumes a MapLibre style~\ac{JSON} document with optional tile statistics and emits a visually equivalent style with optimized tile data.
We formalise the problem and its visual-equivalence constraint (\cref{sec:problem_statement}), scope the work to MapLibre GL~JS (\cref{sec:scope}), and then walk level by level through the three-level pipeline (\cref{sec:approach}) and the passes that implement it.

\section{Formalised Problem Statement}\label{sec:problem_statement}

With the foundations of \cref{ch:theoretical_background} in place, we state the optimization problem formally.

Let $\mathcal{S}$ be the set of syntactically valid MapLibre style~\ac{JSON} documents.\newline
Let $\mathcal{T}$ be the set of tile sets, each a finite family of tiles indexed by source and $(x, y, z)$.\newline
Let $\mathcal{V}$ be the set of \emph{viewing states}: tuples $(\text{lat}_\text{center}, \text{lat}_\text{center}, z, \text{roll}, \text{pitch}, \text{bearing}, \text{viewport size})$, the six degrees of freedom plus the viewport.\newline
Let $\mathcal{I}$ be the set of rasterised images and let $\bot \notin \mathcal{I}$ denote a runtime evaluation error that aborts rendering, as introduced for ill-formed feature-state and global-state references in \cref{sub:feature_state}.
Let MapLibre's deterministic rendering function be
\begin{equation}\label{eq:render_function}
    R : \mathcal{S} \times \mathcal{T} \times \mathcal{V} \to \mathcal{I} \cup \{\bot\}
\end{equation}

The optimizer is a function $\mathcal{O}: \mathcal{S} \times \mathcal{T} \to \mathcal{S} \times \mathcal{T}$ with $(S', T') = \mathcal{O}(S, T)$.
$\mathcal{O}$ pays a one-off preprocessing cost $C_\text{total}(S, T)$, analytically decomposed and bounded in \cref{sub:preprocessing_cost_model}, in order to improve the deployment metrics of \cref{sub:metrics} (load time, heap memory, FPS, with finer-grained diagnostics) for every subsequent rendering session served from $(S', T')$.
The deployment cost is intentionally not formalised as a single closed-form function: it depends on hardware, camera scenario, and tile cache state, and is therefore measured empirically rather than declared analytically.
$\mathcal{O}$ is correspondingly not posed as a single-objective minimisation.
Several passes (e.g.\ selectivity reordering, \cref{sub:expression_optimizations}) deliberately trade a small regression in load time for a larger improvement in FPS, and the encoder competition (\cref{sub:encoding_strategies}) locally minimises encoded byte size as a contributing quantity to load time without committing to any weighting between the deployment axes.
We characterise each pass empirically against the metric battery of \cref{sub:metrics} in the ablation of \cref{ch:experiments}, and the resulting per-axis breakdown is what lets an operator compose the pass set that suits the deployment priorities.

$\mathcal{O}$ must satisfy two formal constraints.
The first constraint is \emph{visual equivalence}.
Writing $\equiv$ for pixel-identity on $\mathcal{I}$,
\begin{equation}\label{eq:visual_equivalence}
    \forall v \in \mathcal{V}:\; R(S, T, v) \neq \bot \;\Rightarrow\; R(S', T', v) \equiv R(S, T, v).
\end{equation}
The constraint is asymmetric in $\bot$.
An optimisation pass may suppress a runtime evaluation error that the original style would have raised, but it must never introduce one that the original would not have raised.
In practice we sample $v$ on a discrete grid of zooms and viewports, and check $\equiv$ by direct image comparison rather than by formal proof.
\Cref{eq:visual_equivalence} doubles as the \emph{soundness criterion} for every pass in this chapter.
Any rewrite whose correctness depends on a MapLibre runtime invariant must either preserve the invariant or reject the rewrite when the invariant cannot be established.

For example, the stats-driven folding pass cannot soundly remove a \mlop{match} arm without a full census of the tile data, because the arm may be reached by features the statistics have not observed.
The pass is consequently gated on a sample rate of~1.0.
Similarly, the metadata-refinement pass must reason about overzoom (\cref{sec:tiles}) when tightening \mlprop{maxzoom}, because an overzoomed tile counts as a viewing state in \cref{eq:visual_equivalence}.

The second constraint is \emph{strict idempotency}:
\begin{equation}\label{eq:idempotency}
    \quad \forall (s, t) \in \mathcal{S} \times \mathcal{T}: \mathcal{O}(\mathcal{O}(s, t)) = \mathcal{O}(s, t)
\end{equation}
Feeding $\mathcal{O}$'s output back as input yields the same pair, ruling out unbounded chained improvement by construction.
The expression-pass fixpoint of \cref{sub:fixpoint} enforces idempotency on the style half.
The deterministic, statistics-driven encoding choices of \cref{sub:data_optimizations} enforce it on the tile half.
Both halves are exercised on the style level by the round-trip proptest and on the tile level by the proptest+libFuzzer harness of \cref{par:encoder_validation}.

\section{Scope}\label{sec:scope}

We target MapLibre GL~JS as the system under test.
As discussed in \cref{sec:styles}, performance research on interactive vector maps is concentrated on Mapbox and MapLibre styles.
The Mapbox codebase is proprietary and may not be evaluated without explicit consent, leaving MapLibre as the only viable open research target.
Within MapLibre, the web renderer \texttt{maplibre-gl-js} is the reference implementation that first tracks new specification features and exposes the richest instrumentation hooks for Puppeteer-driven benchmarking, which is what makes it our benchmark target rather than the Native or mobile renderers that share the same styling engine.

In scope are the three pass families that make up the optimizer pipeline (\cref{sub:pipeline}): expression-level rewrites on style~\ac{JSON} (constant folding, algebraic simplification, statistics-driven folding, selectivity reordering), structural rewrites over a typed style model (dead-layer elimination, metadata refinement, layer merging), and data-level rewrites on the served tile payload (style-driven tile shaving and \ac{MLT} encoding-strategy selection).
Style and data passes are co-designed so that data rewrites are advised by the optimized style, and every pass is bound by the visual-equivalence constraint of \cref{eq:visual_equivalence} as its soundness criterion.

The following concerns are deliberately out of scope:

\begin{description}
    \item[\emph{Energy measurement}] Energy effects are discussed only qualitatively.
We report no \ac{RAPL} or battery deltas because the benchmark host is a desktop workstation rather than a mobile device (\cref{sub:metrics}).
    \item[\emph{Lossy geometry transformations}] Coordinate-precision reduction, Visvalingam-Whyatt, Douglas-Peucker, and polygon coalescing would violate the pixel-identity contract of \cref{eq:visual_equivalence}.
    Writing a visually non-lossy gepgraphic simplification is deferred to future work.
    We therefore quantify the size reductions achievable \emph{without} touching geometry, understating that there is headroom that composition with an geographic simplifier could reach at the cost of changing the style output.
    \item[\emph{Symbol-layer merging}] MapLibre evaluates symbol placement per layer through a tile-wide\newline
    \noindent\texttt{CollisionIndex}, so merging symbol layers would alter label suppression order.
The detailed argument is given along the structural passes (\cref{sub:structural_optimizations}).
    \item[\emph{Proprietary style corpora}] Mapbox Streets, Esri vector basemaps, and similar commercial styles are excluded by their terms of service.
    \item[\emph{Tile schemas other than \ac{OMT}}] Non-\ac{OMT} vector tile sources are excluded from the deep-corpus evaluation, and the stats-driven thresholds used by the pipeline are tuned on Germany~\ac{OMT} only (\cref{sec:gathering_sample}).
    \item[\emph{Online and incremental serving}] The optimizer is an offline, one-shot rebuild over a (style, tile~set) pair.
No integration with a live tile server is assumed, and the optimized artefacts can be served by any standards-compliant tile infrastructure.
    \item[\emph{The reference \ac{MLT} encoder}] The legacy self-contained \ac{MLT} encoder of Tremmel~\&~Zink~\cite{tremmelMapLibreTileNext2025b} is retained as a baseline for the encoding experiments of \cref{ch:experiments}.
It is not a target of optimization in this thesis.
\end{description}

\section{Approach and Basic Building Blocks}\label{sec:approach}

To evaluate the effectiveness of the optimizations strategies developed in this work, we adopted a benchmarking-driven approach.
We first established a reliable performance baseline and then ran cumulative ablation experiments that isolate the contribution of each optimizations pass.
Under the soundness contract of \cref{sec:scope}, the passes from~\cref{sub:pipeline} aim to improve the deployment metrics of \cref{sub:metrics} in exchange for the preprocessing cost $C_\text{total}$ of \cref{sub:preprocessing_cost_model}.
The load-time / FPS trade-off underlying these metrics is discussed in \cref{sec:problem_statement}.

\subsection{Three-Level Optimization Pipeline}\label{sub:pipeline}

\Cref{fig:system_architecture} sketches the optimizer system as a whole.
At the centre sits the \emph{optimiser}, a pipeline of passes at three levels of abstraction.
Around the optimiser are three supporting components, each detailed in its own section.
A code-generation front end (\cref{sub:codegen}) produces Rust types from the upstream MapLibre style specification.
A tile-statistics scanner (\cref{sub:tile_statistics}) profiles \ac{MVT} tiles for the stats-driven and data-level passes.
A validation harness (pixel-exact render tests plus a curated style corpus) exercises both pipelines against the visual-equivalence guarantee of \cref{eq:visual_equivalence}.
The three pass levels at the centre are:

\begin{enumerate}
    \item \textbf{Expression passes} (\cref{sub:expression_optimizations}) operate directly on the \ac{JSON} representation of style expressions, applying peephole rewrites, constant folding, and algebraic simplifications (\cref{sec:compiler_optimizations}) without deserializing the style into typed structures.
    They run to a fixpoint with a hard cap of eight iterations.
    In practice every style in the benchmarking and Maputnik\footnote{\url{https://maputnik.github.io/}, last accessed 2026-05-23} corpora converges within two to three iterations, well below the cap, because human-authored styles tend to reuse a small repertoire of patterns.
    The formal termination measure and confluence (\cref{sub:fixpoint}) discussion are given alongside the pass descriptions.
    \item \textbf{Structural passes} (\cref{sub:structural_optimizations}) deserialize the style into the typed Rust spec produced by the codegen front end and apply transformations that require cross-layer or cross-source reasoning, namely dead layer elimination, metadata refinement, layer merging, and cleanup.
    A cross-phase \emph{filter-to-property constant propagation}, which substitutes equality constraints extracted from a layer's filter into that layer's paint and layout expressions, bridges the expression and structural phases and is detailed in \cref{sub:structural_optimizations}.
    \item \textbf{Data passes} (\cref{sub:data_optimizations}) transform the served tile data via a style-driven pruning advisory and the conversion from \ac{MVT} to \ac{MLT}.
\end{enumerate}

\noindent We chose this ordering to address the \emph{phase ordering problem}~\cite{clickCombiningAnalysesCombining1995}.
Expression passes run first, allowing subsequent structural passes to operate on simplified expressions.
Dead elimination then removes unused layers before layer merging, creating additional merging opportunities.
Finally, the data phase processes the fully optimized style so that pruning recommendations are based on the final set of live layers and properties.

\begin{figure}[htbp]
    \centering
    \begin{tikzpicture}[
        node distance=0.8cm and 1.4cm,
        box/.style={rectangle, draw, rounded corners=2pt, minimum height=0.7cm,
                    minimum width=2.0cm, font=\small, align=center},
        src/.style={box, fill=roleSrc!20},
        mir/.style={box, fill=roleMir!30},
        gen/.style={box, fill=roleGen!20},
        data/.style={box, fill=roleData!25},
        opt/.style={box, fill=roleOpt!35, minimum width=2.4cm},
        test/.style={box, fill=roleTest!25},
        arrow/.style={-Stealth, thick},
        query/.style={-Stealth, thick, densely dashed},
        label/.style={font=\scriptsize\itshape, text=darkgray},
        validationlabel/.style={font=\scriptsize\itshape, text=cvdvermillion},
        validation/.style={
            dashed,
            cvdvermillion,
            ->,
            every edge node/.append style={text=cvdvermillion}
        },
    ]
        % Codegen pipeline (top row)
        \node[src]  (v8)     {\texttt{v8.json}};
        \node[mir,  right=of v8]   (parsed) {ParsedSpec};
        \node[mir,  right=of parsed] (mir)    {MIR};
        \node[gen,  right=of mir]  (serial) {SerialSpec\\{\scriptsize (generated Code)}};

        \draw[arrow] (v8)     -- node[above, label] {decode} (parsed);
        \draw[arrow] (parsed) -- node[above, label] {normalize} (mir);
        \draw[arrow] (mir)    -- node[above, label] {generate} (serial);

        % Data pipeline (bottom row)
        \node[src, below=2.5cm of v8] (proto) {Protobuf\\{\scriptsize (\acs{MVT})}};
        \node[data, right=of proto] (mvt) {MVTLayer};
        \node[data, right=of mvt]   (scan) {DataScanner};

        \draw[query] (proto) -- node[above, label] {decode} (mvt);
        \draw[query] (mvt)   -- node[above, label] {scan tiles} (scan);

        % tangled nodes
        \node[opt, below=2.5cm of serial] (opt) {Optimiser};
        \node[test] (test)  at ($(serial)!0.5!(opt) + (3cm,0)$)  {Render tests};
        \node[test, below=1cm of mir] (corpus) {Style Corpus};

        % tanlged arrows
        \draw[<->,double,thick,densely dashed] (opt) -- node[left,black] {\scriptsize\itshape optimize} (serial);
        \draw[query] (scan) -- node[above, label] {statistics} (opt);
        \draw[validation] (corpus) -- node[left, validationlabel] {validate} (serial);
        \draw[validation] (corpus) -- node[left, validationlabel] {eval} (opt);
        \draw[validation] (test) -- node[right, validationlabel] {validate} (serial);
        \draw[validation] (test) -- node[right, validationlabel] {validate} (opt);
    \end{tikzpicture}
    \caption{
System architecture.
    The \emph{codegen pipeline} (top) transforms the upstream \texttt{v8.json} style specification through three stages into generated Rust types (\texttt{SerialSpec}) used for typed style manipulation, with the MIR providing spec metadata to the optimizer.
    The \emph{data pipeline} (bottom) decodes \acs{MVT} tiles and scans them to produce tile statistics that feed into data-driven optimizations passes.
Render tests validate round-trip correctness and the style corpus is used to both validate the codegen and evaluate the optimizer.}
    \label{fig:system_architecture}
\end{figure}

\subsection{Code Generation from v8.json}\label{sub:codegen}

We design structural passes to operate on Rust types \emph{generated} from the upstream MapLibre GL~JS style specification rather than raw \ac{JSON}.
This shifts spec drift from a manual synchronisation burden to a regenerator invocation, and gives every pass typed access to defaults, value types, and expression capabilities.
Hand-written Rust types are not a viable alternative: MapLibre publishes the specification in its own \ac{JSON}-based \ac{DSL}, and a hand-maintained shadow would require manual synchronisation with every specification change and risk silent divergence.

The generation pipeline (top row of \cref{fig:system_architecture}) proceeds in three stages.
First, it \emph{decodes} the published \texttt{v8.json} into an internal representation.
Second, it \emph{normalizes} the result into a stable intermediate representation (the \texttt{MIR}\footnote{The \texttt{MIR} draws on the normalized-IR design of MLIR~\cite{lattnerMLIRScalingCompiler2021}, providing a stable form between upstream sources and downstream consumers (the codegen emitter and the structural passes).
The dialect, lowering, and pass-manager machinery of the full MLIR framework is not carried over.
The optimizer operates at a single abstraction level with no code-generation target, and round-trips its output to JSON for unmodified MapLibre renderers.}) that resolves cross-cutting concerns.
In particular, it collapses paint and layout properties that the spec defines across multiple sections into per-layer field definitions carrying type, default, and expression-capability metadata.
Third, a code emitter \emph{generates} Rust sources with round-trip serialization support.

This metadata is the substrate for two structural passes.
Layer merging queries the MIR to decide whether a differing paint property supports feature-driven expressions (and can therefore be wrapped in a \mlop{case}/ \mlop{match} over the filter discriminator) or is camera-only and must remain uniform across merged layers.
Default stripping compares property values against the specification defaults encoded in the generated types.

\subsection{Implementation}\label{sub:implementation}

The style optimizer is a standalone Rust binary.
The tile re-encoder described in \cref{sub:data_optimizations} shares the same Cargo workspace and is split into a core encoder/decoder crate and a Java wrapper crate so that a \ac{JVM}-based tile-generation pipeline can consume the encoder without re-implementing it.
The cross-language surface is generated by \emph{Diplomat}\footnote{\url{https://github.com/rust-diplomat/diplomat}, last accessed 2026-05-11}, which produces idiomatic, language-specific bindings directly from annotated Rust types, replacing the earlier \texttt{cbindgen}\footnote{\url{https://github.com/mozilla/cbindgen}, last accessed 2026-05-23}/\texttt{jextract}\footnote{\url{https://github.com/openjdk/jextract}, last accessed 2026-05-23} approach.
The Java wrapper exposes the generated bindings through the Panama \ac{FFM} \ac{API} (JEP~454)\footnote{\url{https://openjdk.org/jeps/454}, last accessed 2026-05-23} rather than conventional \ac{JNI} glue, eliminating hand-written C stubs and giving the \ac{JVM} direct, zero-copy access to off-heap buffers.
This is the boundary shape a columnar encoder requires: each tile is handed over as a batch of parallel \texttt{MemorySegment}s rather than one call per feature, amortising the per-boundary cost over thousands of features in one encoder invocation.

\section{Building Optimizations}\label{sec:building_optimizations}

The implemented passes target three areas of the rendering pipeline: style expressions, structural properties, and tile data.
We classify each pass by \emph{where} it operates and its \emph{data requirements} (none, sampling, or full scan).
Many are analogues of well-known compiler and query-processing techniques.

Each level discharges the visual-equivalence contract of \cref{sec:scope} in a different way.
For expression-level passes, the termination measure $\mu(S)$ (\cref{eq:termination_measure}) witnesses this.
Each rewrite strictly reduces $\mu$ while preserving expression semantics.
For structural passes, we argue soundness per-pass in the paragraphs below.
Dead elimination preserves equivalence because it only removes layers that produce no visible output.
Metadata refinement preserves equivalence because it only tightens bounds that the renderer would enforce anyway (with the overzoom asymmetry detailed below).
Layer merging preserves equivalence via the three invariants of \cref{sub:structural_optimizations}.
For data-level passes, the pruning advisory removes only data that no live style layer references, so the rendering function $R$ receives the same effective input.

\subsection{Expression-Level Optimizations}\label{sub:expression_optimizations}

Expression passes operate directly on the \ac{JSON} representation, applying rewrites that reduce expression complexity and eliminate redundant computation.
\cref{tab:expression_passes} summarizes the implemented passes.

\begin{table}[htbp]
\centering
\begin{tblr}{
  colspec = {p{3.8cm} p{1.8cm} p{9.5cm}},
  rows = {font=\small},
  row{1} = {font=\small\bfseries},
}
\toprule
Pass & Data & Description [compiler analogy] \\
\midrule
Unary simplification         & none     & Unwrap single-argument \mlop{any}/ \mlop{all} wrappers [identity elimination] \\
Expression kind normalization & none     & Rewrite negated comparisons to their direct form, e.g.\ \texttt{[\mlop{"!"}, [\mlop{"=="},a,b]]} $\to$ \texttt{[\mlop{"!="},a,b]} [operator strength reduction] \\
Constant folding              & none     & Evaluate statically determinable expressions: boolean algebra, dead branch elimination, redundant coercion removal [constant folding, \acs{SCCP}] \\
Stats-driven constant folding & sampling & Fold \mlop{has}/ \mlop{get}/ \mlop{geometry-type} expressions using tile statistics; prune unreachable \mlop{match}/ \mlop{in} arms; fold comparisons against known value ranges [sparse conditional constant propagation] \\
Expression simplification     & none     & De~Morgan's laws, boolean flattening, coalesce flattening, \mlop{case}$\to$\mlop{match} rewriting, \mlop{let}/ \mlop{var} inlining, \mlop{any}$\to$\mlop{in} merging, interpolate/step collapsing [algebraic simplification] \\
Default stripping             & none     & Remove paint/layout properties whose value equals the specification default [dead store elimination] \\
Colour minification           & none     & Compress color literals to their shortest representation [literal compression] \\
Selectivity reordering        & sampling & Reorder filter predicates by selectivity so the most discriminating condition is evaluated first [predicate reordering] \\
\bottomrule
\end{tblr}
\caption{
Expression-level optimization passes, each annotated with its data dependency (\emph{none} = static only; \emph{sampling} = requires a tile-statistics scan) and its compiler-optimization analogue.  Passes requiring no statistics run unconditionally.
Sampling passes additionally fold expressions using observed value distributions.}
\label{tab:expression_passes}
\end{table}

Among these passes, we expect constant folding to dominate because reducing an expression to a literal eliminates all per-feature evaluation cost.
The other passes primarily shorten the evaluation path rather than eliminating it.

We organize the expression pass engine as a table of rewrite rules (\cref{tab:all_rewrite_rules}), each gated by a \emph{pass flag} (which pass must be enabled) and scoped to a \emph{rule scope} (whether the rule applies to filters only, to paint/layout properties, or to both).
Selectivity reordering runs as a separate pass after the rewrite table has stabilised.
We further classify rules by their \emph{type}:
	\emph{pure} rules operate on the expression array alone,
	\emph{MIR-aware} rules additionally consult the specification's type metadata (e.g.\ to determine whether an operator is pure or to look up default values), and
	\emph{stats-driven} rules are recursive tree folds that receive tile statistics and layer-level context.
A post-order visitor walks every expression node in the style.
At each node, it tests all applicable peephole rules and the first matching rule fires.
Stats-driven rules run as separate tree folds after the peephole pass, because they require recursive context threading that the node-local peephole architecture does not provide.
We modeled this architecture on peephole optimizers in production compilers.
Individual rules are simple and local, but their composition through fixpoint iteration (\cref{sub:fixpoint}) achieves global simplifications.

\paragraph{Constant Folding}
The constant folding pass implements 17~rewrite rules that collectively perform the role of constant propagation and dead code elimination in a compiler (\cref{sec:compiler_optimizations}).
Boolean algebra rules evaluate \mlop{all}/ \mlop{any} nodes with known operands: a vacuous \texttt{[\mlop{"all"}]} rewrites to \mlval{true} (vacuous truth), a vacuous \texttt{[\mlop{"any"}]} to \mlval{false} (vacuous disjunction), and boolean absorption removes known-true operands from \mlop{all} and known-false operands from \mlop{any}.
Absorption is sound because MapLibre filter operators are pure (the filter language exposes no feature-state mutation, see \cref{sub:feature_state}), so dropping an operand cannot discard an observable side effect, and the asymmetric form of \cref{eq:visual_equivalence} permits dropping an operand whose evaluation would have raised $\bot$.
Comparison rules evaluate binary comparisons with two literal operands.
Dead-branch elimination removes arms from \mlop{case} and \mlop{match} expressions whose guard is statically false.

Contradiction detection identifies filters that are logically unsatisfiable by inspecting siblings within an \mlop{all} node.
When two predicates constrain the same property to disjoint values, e.g.

\begin{lstlisting}[
	style=maplibrestyle,
	basicstyle=\ttfamily\scriptsize,
	frame=single,
	aboveskip=5pt,
	belowskip=5pt
	]
["all",
	["==", ["get", "class"], "a"],
	["==", ["get", "class"], "b"]
]
\end{lstlisting}

\noindent the enclosing \mlop{all} is rewritten to \mlval{false}, which the structural dead elimination pass will later exploit to remove the enclosing layer.
This is the expression-level analogue of unsatisfiable-query detection in a database query planner.

Range tightening detects overlapping range predicates on the same property and replaces them with the tighter bound, e.g.\ \texttt{[\mlop{"all"},~[\mlop{"<"},~x,~\mlval{30}],~[\mlop{"<"},~x,~\mlval{20}]]} becomes \texttt{[\mlop{"<"},~x,~\mlval{20}]}.
Predicate subsumption generalizes this to remove predicates that are logically implied by a sibling.
Pure operator evaluation computes the result of any operator whose arguments are all literals, covering mathematical functions, string operations, and color constructors.

\paragraph{Statistics-Driven Constant Folding}
When tile statistics are available, the optimizer extends constant folding with nine rules that exploit knowledge of the actual data.
These rules receive tile statistics for the relevant source layer and fold expressions that are statically indeterminate but data-determinate:

\begin{description}
    \item[\textbf{Property presence folding}] if a property is present in every feature, \texttt{[\mlop{"has"},~\mlval{"p"}]} folds to \mlval{true}.
If present in none, it folds to \mlval{false}.
    \item[\textbf{Single-value folding}] if a property has cardinality~1 and is present in all features, \texttt{[\mlop{"get"},~\mlval{"p"}]} folds to the single known value.
    \item[\textbf{Geometry-type folding}] if a source layer contains only one geometry type, comparisons against \mlop{geometry-type} fold to a boolean.
    \item[\textbf{Range folding}] numeric comparisons against a property with known data range $[\mathrm{min},\, \mathrm{max}]$ are folded when the predicate is trivially satisfied or violated.
    We apply the following folding rules:
    \begin{equation}\label{eq:range_folding}
    \begin{aligned}
        p < n  &\;\to\; \mathtt{false} \;\text{if}\; \mathrm{min} \geq n, &
        p < n  &\;\to\; \mathtt{true}  \;\text{if}\; \mathrm{max} < n, \\
        p \leq n &\;\to\; \mathtt{false} \;\text{if}\; \mathrm{min} > n, &
        p \leq n &\;\to\; \mathtt{true}  \;\text{if}\; \mathrm{max} \leq n, \\
        p > n  &\;\to\; \mathtt{false} \;\text{if}\; \mathrm{max} \leq n, &
        p > n  &\;\to\; \mathtt{true}  \;\text{if}\; \mathrm{min} > n, \\
        p \geq n &\;\to\; \mathtt{false} \;\text{if}\; \mathrm{max} < n, &
        p \geq n &\;\to\; \mathtt{true}  \;\text{if}\; \mathrm{min} \geq n, \\
        p = n  &\;\to\; \mathtt{false} \;\text{if}\; n < \mathrm{min} \lor n > \mathrm{max}, &
        p \neq n &\;\to\; \mathtt{true}  \;\text{if}\; n < \mathrm{min} \lor n > \mathrm{max}.
    \end{aligned}
    \end{equation}
    All cases additionally require that the property is present in all features (otherwise the comparison is undefined for absent features).
    We apply the equality and disequality rules only to integer-typed columns, because floating-point equality is not a reliable oracle: a stats min/max pair derived from observed doubles may legitimately exclude a literal by a trailing-bit margin, but the filter is still defined (and, for correctness, treated as non-foldable).
    We apply strict-inequality folds on floating-point columns at the reported min/max values exactly.
Stats collection uses the observed extrema rather than synthetic bounds, so the boundary comparison coincides with a value that actually occurs, which is the only case where the fold is unconditionally sound.
    \item[\textbf{Match/in arm pruning}] the pass removes arms referencing values absent from the tile statistics.
If all arms are pruned, the expression collapses to its fallback.
    \item[\textbf{Match arm reordering}] the pass reorders arms by descending frequency so the renderer evaluates the most common case first.
    \item[\textbf{Data-ramp pruning}] for \mlop{step} and \mlop{interpolate} expressions driven by a numeric property with known data range $[\mathrm{min},\, \mathrm{max}]$, the pass prunes stops outside the observed range.
    In a \mlop{step} expression, it removes stops with thresholds above $\mathrm{max}$.
When the property is present in all features, it absorbs stops below $\mathrm{min}$ into the default.
    In an \mlop{interpolate} expression, it prunes stops outside the $[\mathrm{min},\, \mathrm{max}]$ interval symmetrically.
\end{description}

\noindent A cardinality threshold of \texttt{200} governs the granularity of value tracking: when the number of distinct values exceeds this threshold, the statistics switch from full value counts to cardinality-only mode.
This bounds the cost on high-cardinality properties (names, identifiers) while retaining full information for the low-cardinality properties (class, subclass, administrative level) that benefit most from data-driven folding.
We motivated the choice of 200 with the OpenMapTiles schema: every categorical property that is used as a \mlop{match}/ \mlop{case} discriminator in the benchmark styles (\texttt{\mlval{class}}, \texttt{\mlval{subclass}}, \texttt{\mlval{admin\_level}}, \texttt{\mlval{kind}}, land-use codes) has fewer than 200 distinct values across the Germany tileset, whereas high-cardinality string columns (feature names, identifiers) exceed it by one or more orders of magnitude.
The threshold therefore separates the \enquote{folding-useful} properties from the \enquote{folding-useless} ones without tuning for any specific style.

We gate folds whose effect is irreversible (those that change the set of rendered features) on a full-census sample rate, as discussed in \cref{sub:tile_statistics}.
At lower sample rates the statistics are used only for reordering and selectivity estimation.

\paragraph{Expression Simplification}
Applying the algebraic identities from \cref{sec:compiler_optimizations}, De~Morgan's laws push negations inward to expose further folding:
%
\begin{quote}
\texttt{[\mlop{"!"},~[\mlop{"all"},~a,~b]]}~$\to$~\texttt{[\mlop{"any"},~[\mlop{"!"},~a],~[\mlop{"!"},~b]]}
\end{quote}
%
\noindent Boolean flattening merges nested \mlop{all}/ \mlop{any}:
%
\begin{quote}
\texttt{[\mlop{"all"},~[\mlop{"all"},~a,~b],~c]}~~$\to$~~\texttt{[\mlop{"all"},~a,~b,~c]}
\end{quote}
%
\noindent The \mlop{any}$\to$\mlop{in} rewrite detects two or more equality tests against the same property within an \mlop{any} and merges them into a single \mlop{in} expression that the renderer can evaluate as a set membership test.
Within \mlop{match} expressions, the pass merges duplicate arms and removes arms whose output equals the fallback.

\begin{sloppypar}
The pass detects boolean \mlop{match} patterns (\texttt{[\mlop{"match"},\allowbreak~x,\allowbreak~labels,\allowbreak~\mlval{true},\allowbreak~\mlval{false}]}) and rewrites them to \texttt{[\mlop{"in"},\allowbreak~x,\allowbreak~[\mlop{"literal"},\allowbreak~labels]]}.
Interpolation and step ramp simplification detects ramps, where all stops map to the same value and collapses them to that constant.
\mlop{let}/ \mlop{var} inlining replaces singly-used variable references with the bound expression, eliminating the binding overhead.
\end{sloppypar}

\paragraph{Default Stripping}
The pass removes properties whose value equals the specification default.
We derive the defaults from the code-generated Rust types: each property's default is known from the upstream \texttt{v8.json}.
The pass strips only plain scalar values.
It never removes expression-valued properties, even if they would evaluate to the default at all zoom levels, to avoid the risk of changing renderer behavior when the expression's mere presence affects layout (e.g.\ \texttt{\mlprop{text-radial-offset}} suppresses \texttt{\mlprop{text-offset}} regardless of its value).

\paragraph{Selectivity Reordering}\label{sec:selectivity_reordering}
For variadic boolean expressions (\mlop{all}, \mlop{any}), the pass reorders predicates by estimated selectivity to exploit short-circuit evaluation, applying the selectivity-estimation framework of \cref{sec:query_optimizations}.
For \mlop{all} (conjunction), it sorts predicates by ascending selectivity (the predicate most likely to be false is evaluated first).
For \mlop{any} (disjunction), it sorts predicates by descending selectivity (the predicate most likely to be true is evaluated first).
The selectivity estimation formulas follow the independence assumption introduced by Selinger et al.~\cite{selingerAccessPathSelection1979}.
Let $N$ denote the total number of features in the source layer, and $\mathrm{freq}(v)$ the number of features where property~$p$ equals~$v$:
\begin{equation}\label{eq:selectivity}
    \begin{aligned}
        \mathrm{sel}(p = v)   &= \frac{\mathrm{freq}(v)}{N}  &
        \mathrm{sel}(p < n)   &= \frac{1}{N}\!\sum_{v < n}\! \mathrm{freq}(v)  \\[4pt]
        \mathrm{sel}(\mathtt{\mlop{geometry\text{-}type}}\; t) &= \frac{\mathrm{count}(t)}{N} &
        \mathrm{sel}(p \in V)  &= \min\!\Bigl(1,\;\sum_{v \in V}\! \mathrm{sel}(p = v)\Bigr) \\[4pt]
        \mathrm{sel}(\mathtt{\mlop{has}}\; p) &= \frac{\mathrm{present}(p)}{N} &
        \mathrm{sel}(p \neq v) &= 1 - \mathrm{sel}(p = v) \\[4pt]
    \end{aligned}
\end{equation}
We estimate compound predicates under the standard independence assumption:
\begin{equation}\label{eq:selectivity_compound}
    \begin{aligned}
        \mathrm{sel}(\mlop{all} \quad p_0 \ldots p_i) &= \prod_i \mathrm{sel}(p_i) \\[4pt]
        \mathrm{sel}(\mlop{any} \quad p_0 \ldots p_i) &= 1 - \prod_i \bigl(1 - \mathrm{sel}(p_i)\bigr) \\[4pt]
    \end{aligned}
\end{equation}

\noindent The independence assumption is known to be imprecise for correlated predicates~\cite{selingerAccessPathSelection1979}, but as discussed in \cref{sec:query_optimizations} only a reordering heuristic consumes the selectivity signal here, not a cost-based plan chooser: what matters is the \emph{ordering} of predicates by selectivity, not their absolute values.
Correlation between map-style predicates (e.g.\ \texttt{\mlval{class}} and \texttt{\mlval{subclass}} in OpenMapTiles) distorts the absolute selectivities but rarely inverts their relative ordering, so the cheap estimator suffices.
Because the rendering benefit of predicate reordering depends on the proportion of evaluation cost attributable to predicate ordering (expected to be small for styles with short filters) we include this pass as an exploratory extension and do not separately benchmark its isolated impact.

\paragraph{Termination and Confluence of the Fixpoint}
The expression-level phase applies the entire rewrite rule set until the style \ac{JSON} stops changing.
In the \ac{TRS} framing of \cref{sub:fixpoint}, this is the standard normal-form criterion.
We impose a hard cap of \texttt{8} iterations.
This is the \emph{formal} termination guarantee.
The fixpoint loop executes at most 8 passes regardless of rule behavior.

The majority of the rules described above strictly decrease a lexicographic measure
\begin{equation}\label{eq:termination_measure}
    \mu(S) = (n_{\mathrm{AST}}(S),\, d_{\max}(S),\, n_{\mathrm{lit}}(S))
\end{equation}
where $n_{\mathrm{AST}}$ is the total number of \ac{AST} nodes in all expressions of $S$, $d_{\max}$ is the maximum expression nesting depth, and $n_{\mathrm{lit}}$ is the number of non-literal nodes.
Rules either reduce $n_{\mathrm{AST}}$ (e.g.\ constant folding, boolean absorption, dead-branch elimination), keep $n_{\mathrm{AST}}$ unchanged but strictly reduce $d_{\max}$ (e.g.\ boolean flattening), or keep both unchanged but strictly reduce $n_{\mathrm{lit}}$ (e.g.\ redundant coercion removal).
Two further rules (\mlop{any}$\to$\mlop{in} merging and predicate reordering) are either strictly monotone in the measure (\mlop{any}$\to$\mlop{in} reduces $n_{\mathrm{AST}}$ because the merged \mlop{in} has fewer nodes than the original chain of equalities) or semantically idempotent after one application (selectivity reordering sorts predicates into a canonical order).

Two rules can \emph{temporarily increase} $n_{\mathrm{AST}}$ within a single iteration, which is why the measure alone does not constitute a termination proof.
De~Morgan's law transforms \texttt{[\mlop{"!"}, \allowbreak[\mlop{"any"}, \allowbreak A, \allowbreak B]]} (4~nodes) into \texttt{[\mlop{"all"}, \allowbreak[\mlop{"!"}, \allowbreak A], \allowbreak[\mlop{"!"}, \allowbreak B]]} (5~nodes), and distributive factoring may introduce boolean identity elements in vacuous-remainder cases.
In both cases, subsequent rules within the same or next iteration typically compensate: boolean flattening reduces nesting depth, and constant folding absorbs the identity elements, so the net effect across iterations is a decrease in $\mu$.
In the general case, however, compensating rules are not guaranteed to fire (e.g., when De~Morgan produces negated predicates that admit no further simplification at the top level of a filter).
In practice, all styles in the benchmark corpus converge within three iterations, well below the hard cap, because the monotone majority of rules drives $\mu$ toward zero and the non-monotone rules fire at most once per expression.

The rewrite system is not confluent.
The first matching rule fires (table iteration order defines implicit priority), and selectivity reordering depends on data-dependent statistics.
Different rule orderings could therefore yield different but visually equivalent output.
Visual equivalence (the soundness criterion of~(\ref{eq:visual_equivalence})) is preserved regardless, because every individual rule preserves expression semantics.



\paragraph{Concrete Examples}
To illustrate how expression passes compose, \cref{fig:expr_before_after} shows four representative transformations.

\begin{figure}[ht]
	\centering

	% =========================================================
	% A
	% =========================================================
	\begin{subfigure}[t]{\linewidth}
		\centering

		\begin{minipage}[c]{0.42\linewidth}
			\centering
			\begin{lstlisting}[
				style=maplibrestyle,
				basicstyle=\ttfamily\scriptsize,
				frame=single,
				aboveskip=0pt,
				belowskip=0pt
				]
["any",
	["all",
		["has","name"],
		["==", ["get","class"], "road"]
	],
	["all",
		["has","name"],
		["==", ["get","class"], "rail"]
	]
]
			\end{lstlisting}
		\end{minipage}
		\hfill
		\begin{minipage}[c]{0.08\linewidth}
			\centering
			\vfill
			{\Huge$\Longrightarrow$}
			\vfill
		\end{minipage}
		\hfill
		\begin{minipage}[c]{0.42\linewidth}
			\centering
			\begin{lstlisting}[
				style=maplibrestyle,
				basicstyle=\ttfamily\scriptsize,
				frame=single,
				aboveskip=0pt,
				belowskip=0pt
				]
["all",
	["has", "name"],
	["in",
		["get", "class"],
		["literal", ["road", "rail"]]
	]
]
			\end{lstlisting}
		\end{minipage}

		\caption{
Expression rewriting via distributive factoring and \mlop{any}$\to$\mlop{in} promotion.}
		\label{fig:expr_a}
	\end{subfigure}

	\vspace{4mm}

	% =========================================================
	% B
	% =========================================================
	\begin{subfigure}[t]{\linewidth}
		\centering

		\begin{minipage}[c]{0.42\linewidth}
			\centering
			\begin{lstlisting}[
				style=maplibrestyle,
				basicstyle=\ttfamily\scriptsize,
				frame=single,
				aboveskip=0pt,
				belowskip=0pt
				]
filter: ["==", ["get", "class"], "road"]
paint:
	fill-color: ["match",
		["get", "class"],
			"road", "#333",
			"rail", "#666",
			"#999"
		]
			\end{lstlisting}
		\end{minipage}
		\hfill
		\begin{minipage}[c]{0.08\linewidth}
			\centering
			\vfill
			{\Huge$\Longrightarrow$}
			\vfill
		\end{minipage}
		\hfill
		\begin{minipage}[c]{0.42\linewidth}
			\centering
			\begin{lstlisting}[
				style=maplibrestyle,
				basicstyle=\ttfamily\scriptsize,
				frame=single,
				aboveskip=0pt,
				belowskip=0pt
				]
filter: ["==", ["get", "class"], "road"]
paint:
	fill-color: "#333"
			\end{lstlisting}
		\end{minipage}

		\caption{
Filter to Property Constant Propagation simplifying data-driven to constant.}
		\label{fig:expr_b}
	\end{subfigure}

	\vspace{4mm}

	% =========================================================
	% C
	% =========================================================
	\begin{subfigure}[t]{\linewidth}
		\centering

		\begin{minipage}[c]{0.42\linewidth}
			\centering
			\begin{lstlisting}[
				style=maplibrestyle,
				basicstyle=\ttfamily\scriptsize,
				frame=single,
				aboveskip=0pt,
				belowskip=0pt
				]
["all",
	["==", ["get", "class"], "road"],
	["==", ["get", "class"], "road"]
]
			\end{lstlisting}
		\end{minipage}
		\hfill
		\begin{minipage}[c]{0.08\linewidth}
			\centering
			\vfill
			{\Huge$\Longrightarrow$}
			\vfill
		\end{minipage}
		\hfill
		\begin{minipage}[c]{0.42\linewidth}
			\centering
			\begin{lstlisting}[
				style=maplibrestyle,
				basicstyle=\ttfamily\scriptsize,
				frame=single,
				aboveskip=0pt,
				belowskip=0pt
				]
["==", ["get", "class"], "road"]
			\end{lstlisting}
		\end{minipage}

		\caption{
Duplicate predicate elimination via subsumption.}
		\label{fig:expr_c}
	\end{subfigure}

	\vspace{4mm}

	% =========================================================
	% D
	% =========================================================
	\begin{subfigure}[t]{\linewidth}
		\centering

		\begin{minipage}[c]{0.42\linewidth}
			\centering
			\begin{lstlisting}[
				style=maplibrestyle,
				basicstyle=\ttfamily\scriptsize,
				frame=single,
				aboveskip=0pt,
				belowskip=0pt
				]
["all",
	["!", ["==", ["get", "class"], "road"]],
	["!", ["==", ["get", "class"], "rail"]]
]
			\end{lstlisting}
		\end{minipage}
		\hfill
		\begin{minipage}[c]{0.08\linewidth}
			\centering
			\vfill
			{\Huge$\Longrightarrow$}
			\vfill
		\end{minipage}
		\hfill
		\begin{minipage}[c]{0.42\linewidth}
			\centering
			\begin{lstlisting}[
				style=maplibrestyle,
				basicstyle=\ttfamily\scriptsize,
				frame=single,
				aboveskip=0pt,
				belowskip=0pt
				]
["all",
	["!=", ["get", "class"], "road"],
	["!=", ["get", "class"], "rail"]
]
			\end{lstlisting}
		\end{minipage}

		\caption{
Negation pushdown and operator normalization.}
		\label{fig:expr_d}
	\end{subfigure}

	\caption{
Examples of expression-level transformations applied to representative style inputs.
Each panel shows a before/after \ac{JSON} example pair demonstrating how these optimisations interact.}
	\label{fig:expr_before_after}
\end{figure}

\clearpage
\subsection{Structural Optimizations}\label{sub:structural_optimizations}

Structural passes operate on typed Rust representations of the full style document, enabling cross-layer and cross-source reasoning.
\cref{tab:structural_passes} summarizes the implemented passes.

\begin{table}[htbp]
\centering
\begin{tblr}{
  colspec = {p{3.8cm} p{1.8cm} p{9.5cm}},
  rows = {font=\small},
  row{1} = {font=\small\bfseries},
}
\toprule
Pass & Data & Description [compiler analogy] \\
\midrule
Metadata stripping           & none      & Remove \texttt{metadata} fields at the style root and layer level [debug info stripping] \\
Dead layer/source elimination & none / sampling & Remove layers with always-false filters, geometry-type mismatches (with stats), or empty source-layers; prune unreferenced sources [dead code elimination] \\
Metadata refinement          & none / sampling & Extract \mlprop{minzoom}/ \mlprop{maxzoom} bounds from filter predicates and paint ramps; tighten bounds using tile coverage statistics; remove consumed zoom predicates from filters [range analysis] \\
Cleanup                      & none      & Remove empty \mlprop{paint}/\mlprop{layout} sections, invisible layers (\texttt{\mlprop{visibility}:~\mlval{none}}, zero opacity), trivially \mlval{true} filters, and root-level properties at specification defaults [dead code elimination] \\
Layer merging                & none      & Merge identical adjacent layers by synthezising \mlop{case}/ \mlop{match} expressions, with zoom-tolerant matching and sort-key preservation [basic block merging] \\
Source zoom tightening       & none      & Tighten source \mlprop{minzoom}/ \mlprop{maxzoom} based on the zoom ranges of layers that reference them [liveness analysis] \\
\bottomrule
\end{tblr}
\caption{
Structural optimization passes operating on the typed Rust style representation, each annotated with its data dependency and compiler-optimization analogue.  Passes marked \emph{none/sampling} run unconditionally but gain additional precision when tile statistics are available.}
\label{tab:structural_passes}
\end{table}


\paragraph{Dead Layer and Source Elimination}
The pass treats a layer as dead when its filter is the literal \mlval{false} (produced by contradiction detection or stats-driven folding in the expression phase), when its source layer carries zero features per tile statistics, or when its geometry type does not match the source layer's available types.
After removing dead layers, the pass prunes sources that no surviving layer references.
This mirrors dead code elimination followed by unreachable function removal in a compiler (\cref{sec:compiler_optimizations}).
Its downstream effect on the tile pruning advisory (\cref{sub:tile_shaving}) is the entry point for the data-level cascade evaluated in \cref{sub:mlt_results}.

Geometry-type mismatch is particularly effective on schemas such as OpenMapTiles, where a single source layer (e.g.\ \texttt{transportation}) carries mixed geometry. \texttt{\mlval{fill}} layers require polygons, \texttt{\mlval{circle}} requires points, \texttt{\mlval{line}} draws on polygons or lines but not points, and \texttt{\mlval{symbol}} layers draw on all three and so are never killed by this rule.
The check requires tile statistics, so we provide both static and stats-driven variants of the pass.

\paragraph{Metadata Refinement}
This pass extracts implicit zoom bounds from filter predicates and paint ramps, and lifts them into the layer's explicit \mlprop{minzoom}/ \mlprop{maxzoom} metadata.
Accurate metadata allows the renderer to skip entire layers at zoom levels where they are invisible, avoiding filter evaluation, partially geometry decoding, and draw-call submission.

\emph{Filter-based extraction} walks the filter expression tree and recognizes zoom predicates of the form \texttt{[\mlop{">="},~[\mlop{"zoom"}],~$z$]} or \texttt{[\mlop{"<="},~[\mlop{"zoom"}],~$z$]}.
Within an \mlop{all} node, it intersects the extracted bounds (the tighter of the bounds applies).
Within an \mlop{any} node, it unions them (the layer is visible if any branch is satisfied).
The pass rounds extracted fractional bounds to integer zoom levels because tile zoom levels are integral.
For a comparison operator with threshold~$n$, the extracted zoom interval is:
\begin{equation}\label{eq:zoom_bounds}
    \begin{aligned}
        \texttt{>=}\; n &\;\Rightarrow\; [n,\;\infty) \\
        \texttt{>}\; n  &\;\Rightarrow\; [\lfloor n \rfloor + 1,\;\infty) \\
        \texttt{<=}\; n &\;\Rightarrow\; [0,\; n]\\
        \texttt{<}\; n  &\;\Rightarrow\; [0,\; \lceil n \rceil - 1] \\
    \end{aligned}
\end{equation}
Within an \mlop{all} node, it intersects extracted intervals:
\[
z^{*} = [l_1, u_1] \land [l_2, u_2] = [\max(l_1, l_2),\; \min(u_1, u_2)]
\]
Within an \mlop{any} node, it unions them:
\[
z^{*} = [l_1, u_1] \lor [l_2, u_2] = [\min(l_1, l_2),\; \max(u_1, u_2)]
\]
After extraction, the pass removes the consumed zoom predicates from the filter, often simplifying it to \mlval{true}.

\emph{Paint-based visibility analysis} infers zoom bounds from paint properties that control visibility.
For opacity properties, the pass examines \mlop{interpolate} and \mlop{step} expressions to find the zoom level at which the property transitions from zero to a non-zero value.
For example, a fill layer whose \mlprop{fill-opacity} is a step function that outputs~0 below zoom~12 and~1 above is effectively invisible below zoom~12.
The pass therefore raises \mlprop{minzoom} to~12.
For symbol layers, \emph{either} text or icon opacity must be non-zero (the layer is visible if any symbol component is visible).
For all other layer types, \emph{both} size and opacity must be non-zero.

\emph{Stats-driven refinement} tightens bounds using actual tile coverage: if a source layer's first tile appears at zoom~4, the pass raises the layer's \mlprop{minzoom} to~4 regardless of what the filter permits, since there is no data to render below that zoom.
This refinement only tightens bounds.
It never loosens them due to \cref{eq:visual_equivalence}.

A subtle but load-bearing correctness concern is MapLibre's \emph{overzoom} behavior: when a client requests a tile at a zoom level beyond a source's \mlprop{maxzoom}, the renderer reuses the nearest available tile from a lower zoom and upscales it.
A naive stats-driven tightening (\enquote{the last tile containing polygons appears at zoom~12, therefore set \mlprop{maxzoom}=\mlval{12}}) would change the rendered output if the layer is visible at zoom~13 or~14 via overzoom from zoom~12.

Because our work requires  \cref{eq:visual_equivalence}, the metadata refinement pass enforces two asymmetric rules:
\begin{enumerate}
    \item \textbf{\mlprop{minzoom} is symmetric.}
    Vector tiles are not \emph{underzoomed} (the renderer will not synthezise a coarser tile from a finer one) so raising \mlprop{minzoom} to the first data-carrying zoom never removes a rendered pixel.
    \item \textbf{\mlprop{maxzoom} is asymmetric.}
    The refinement pass tightens the layer's \mlprop{maxzoom} to the last data-carrying zoom $z^{*}$ only when $z^{*}$ is strictly less than the \emph{source's} \mlprop{maxzoom}.
Otherwise, the layer's tiles at $z^{*}$ will be overzoomed to produce higher zoom levels, and the layer remains rendered beyond $z^{*}$.
\end{enumerate}
The zoom-folding rule that collapses zoom comparisons in filters is similarly one-sided: only \mlprop{minzoom} is used as a lower bound for the fold, because a layer's effective \mlprop{maxzoom} is not an upper bound on the camera zoom once overzooming is considered.
The source-zoom-tightening pass follows the same principle: it tightens a source's \mlprop{maxzoom} only as far as the maximum \mlprop{maxzoom} of any layer referencing it, which already respects the overzoom rule above, and it never tightens below the source's declared bounds.
Taken together, these rules ensure that no optimizations can remove a pixel the original style would have rendered via overzoom.

\paragraph{Cleanup}
The cleanup pass performs four operations: it removes empty \mlprop{paint} and \mlprop{layout} sections that may have become empty after default stripping or stats-driven folding, removes invisible layers (\texttt{\mlprop{visibility}:~\mlval{none}}, or zero opacity that is not an expression), removes trivially true filters (\texttt{[\mlop{"literal"},~\mlval{true}]}) that are semantically equivalent to no filter, and strips root-level style properties that equal their specification defaults (bearing, pitch, roll, transition).
Cleanup runs after metadata refinement because refinement may render filters trivially true or make layers invisible.

\paragraph{Layer Merging}
Layer merging identifies adjacent layers of the same type with identical paint and layout properties (modulo filter differences) and merges them into a single layer.
The motivation is that fewer layers means fewer rendering passes.
The renderer processes layers sequentially, and each layer incurs setup costs (binding textures, configuring blend state, submitting draw calls).
\Cref{lst:painter_loop} shows the two render-pass loops the painter runs over \texttt{layerIds}: the opaque pass walks them top-to-bottom (\texttt{currentLayer} from \texttt{layerIds.length - 1} down to \texttt{0}), the translucent pass walks them bottom-to-top, and each iteration of either loop dispatches one \texttt{this.renderLayer(...)} call.
Each layer eliminated by the merge therefore shortens \texttt{layerIds} by one entry and removes exactly one \texttt{renderLayer} dispatch from each pass.

A second per-layer loop in the renderer has the opposite implication for the merging pass and makes \mlval{symbol} layers non-mergeable. \Cref{lst:placement} shows the per-layer processing of symbol placement, where \stepcirc{1} constructs a fresh \texttt{LayerPlacement} instance when a symbol layer's placement begins, and \stepcirc{2} tears it down before the next layer starts.
The underlying \texttt{Placement} object (and its \texttt{CollisionIndex}) is shared across all layers via \texttt{this.placement}, which is what makes the \texttt{LayerPlacement} scope per-layer rather than global.
Symbols from earlier layers are fully placed before later layers begin, so layer ordering determines placement priority.
Merging two symbol layers into one would interleave their symbols within a single \texttt{LayerPlacement}, changing which symbols take priority and therefore altering collision outcomes.

With both renderer-side loops accounted for, we can describe the algorithm itself.
The merging pass proceeds in three phases, illustrated in \cref{fig:layer_merge_example}.
\emph{Phase}~\stepcirc{1} (candidate identification) scans the layer list for contiguous runs of layers that share the same type, source, and source-layer, and whose differing properties are all mergeable.
A property is mergeable if, according to the MIR field metadata,
  it supports feature-driven expressions (i.e.\ its \texttt{expression.feature} flag is \mlval{true}, meaning it can vary per feature within a layer) and
  it is not a zoom ramp (\mlop{interpolate} or \mlop{step} over \texttt{[\mlop{"zoom"}]}).
  Both of these must remain at the top level of a property expression and cannot be nested inside a \mlop{case}/ \mlop{match} wrapper.
Doing so would introduce an $\bot$ which is not allowed according to \cref{eq:visual_equivalence}.
Camera-only properties (such as \texttt{\mlprop{line-dasharray}} or \texttt{\mlprop{*-translate}}) must be uniform across all layers in a merge group, since they cannot be conditioned on feature attributes.
Layer types that support sort keys (\texttt{\mlval{fill}}, \texttt{\mlval{line}}, \texttt{\mlval{circle}}, \texttt{\mlval{fill-extrusion}}) are candidates.
The \texttt{\mlval{symbol}} type is excluded for the placement-scope reason given above.

The merge pass has to preserve three invariants that collectively witness the visual-equivalence constraint of \cref{sec:problem_statement}.

\textbf{Coverage.}
Every feature drawn by the original $k$ layers must still be drawn by the merged layer with the same paint and layout properties.
This holds because the merged filter is the disjunction $\bigvee_i F_i$, the merged paint/layout expression is a \mlop{case}/\mlop{match} on the same discriminator that selects the original per-layer value, and the optimizer refuses the merge whenever a differing property lacks \texttt{expression.feature} support.

\textbf{Painter order.}
The per-layer draw order must be preserved so the top-most pixel at each screen position is unchanged.
For layer types that expose \texttt{\{\mlprop{fill,line,circle}\}\mlprop{-sort-key}}, the pass synthezises a \mlop{case}/\mlop{match} sort key mapping each original layer's features to an ascending integer in declaration order, which the MapLibre renderer then applies before drawing.
The pass refuses to merge any candidate that already carries an explicit sort key, since combining a user-declared key with a synthezised one would over-determine the ordering.
\texttt{\mlprop{fill-extrusion}} layers are eligible in principle but inherit the same sort-key restriction, and their relative ordering is in any case dominated by the depth buffer rather than the painter's algorithm, so the synthezised key only breaks ties at co-planar surfaces.

\textbf{Layer scope.}
Any MapLibre runtime behavior whose scope is the individual layer (per-layer symbol placement against the tile-wide \texttt{CollisionIndex}, \texttt{\mlop{feature-state}}, \texttt{\mlprop{*-translate}}, and other camera-only properties) must not have that scope widened by the merge.
This is guaranteed by excluding symbol layers from the candidate set, by the \texttt{expression.feature} gate forcing camera-only properties to stay uniform, and by refusing to merge layers with differing non-mergeable properties.

\paragraph{Scope of Filter-to-Property Propagation}
The filter-to-property constant propagation pass extracts equality constraints from a layer's filter and substitutes them into the layer's paint and layout expressions.
Its soundness depends on a MapLibre evaluation-order invariant.
The filter is evaluated first, and paint/layout expressions are only evaluated for features that passed the filter.
\Cref{lst:bucket_populate} shows this invariant in the \texttt{FillBucket.populate} method, where \stepcirc{1} rejects a feature before any paint or layout expression is evaluated.
Any equality constraint extracted from the filter is guaranteed to hold for every feature that reaches \stepcirc{2} and can be safely substituted into the layer's paint/layout expressions.
This holds for the data-driven expression variants used across all MapLibre versions this thesis targets (\mlop{match}, \mlop{case}, \mlop{get}, etc.).

Propagation is strictly \emph{intra-layer} (intra-procedural in compiler terms, \cref{sec:compiler_optimizations}).
Constraints from one layer never flow to a sibling that shares the same source.
Phase~\stepcirc{1} also ignores \texttt{[\mlop{"feature-state"}, \ldots]} references, since feature state is set from the client at render time and cannot be statically constrained by the filter.

\emph{Phase}~\stepcirc{2} (match-pattern detection) checks whether all filters within a candidate group match the pattern \texttt{[\mlop{"=="},~[\mlop{"get"},~P],~\mlop{literal}]} on the same property~$P$.
If so, the merged filter is a compact \mlop{match} expression rather than a verbose \mlop{case} chain, and merged paint/layout properties also use \mlop{match} with the same discriminator.
This pattern is common in practice.
Road styles frequently define separate layers for \texttt{highway} = \{\texttt{motorway}, \texttt{trunk}, \texttt{primary}, \texttt{secondary}, \texttt{tertiary\}.
These layers commonly differ only in width, color and layering, and for casing major roads there are usually two such layers per class.

\Cref{fig:layer_merge_example} works through exactly this case: three adjacent \texttt{\mlval{line}} layers keyed on \texttt{\mlval{class}} collapse into a single layer whose filter, \mlprop{line-color}, \mlprop{line-width}, and synthezised \mlprop{line-sort-key} are all \mlop{match} expressions over the same discriminator, with Phase~\stepcirc{1} (mergeability) and Phase~\stepcirc{3} (identical zoom ranges) discharged trivially.

\emph{Phase}~\stepcirc{3} (zoom-tolerant merging) handles layers whose zoom ranges differ slightly.
It wraps each sub-filter with zoom guards (\texttt{[\mlop{">="},\allowbreak~[\mlop{"zoom"}],\allowbreak~lo]}, \texttt{[\mlop{"<="},\allowbreak~[\mlop{"zoom"}],\allowbreak~hi]}) and unions the zoom ranges into the merged layer's bounds, trading a slightly more complex filter for one fewer layer.

The synthezised \emph{sort key} that preserves inter-layer draw order uses the discriminator value as its index under a \mlop{match} pattern and tests filters in reverse order (last painter wins) under a \mlop{case} pattern.
The pass emits a sort key only for layer types that support \texttt{\mlprop{*-sort-key}} and that do not already define one.

\paragraph{Source Zoom Tightening}
After all layer-level passes, this pass tightens source \mlprop{minzoom}/ \mlprop{maxzoom} bounds to the effective zoom range of the layers that reference each source.
This is a liveness analysis in the sense of \cref{sec:compiler_optimizations}.
A source is \enquote{live} at a zoom level only if at least one layer referencing it is active at that zoom.
Tighter source bounds allow the tile server (or client-side tile manager) to skip tile requests for zoom levels where no layer needs data from that source.

\begin{figure}[H]
    \centering
    \resizebox{\linewidth}{!}{\hfuzz=50pt\begin{tikzpicture}[
        node distance=0.35cm,
        layer/.style={rectangle, draw, rounded corners=2pt, minimum width=3.2cm,
                      minimum height=0.55cm, font=\scriptsize, align=left},
        merged/.style={rectangle, draw, rounded corners=2pt, minimum width=4.8cm,
                       font=\scriptsize, align=left, inner sep=4pt},
        phase/.style={rectangle, draw, dashed, rounded corners=3pt, fill=gray!8,
                      inner sep=5pt},
        arrow/.style={-Stealth, thick},
        label/.style={font=\scriptsize\itshape, text=black!90},
        check/.style={font=\scriptsize, text=black!90},
    ]
        % Input layers (left)
        \node[layer]   (l0) {\mlval{line}\\ \texttt{\mlprop{filter}: [\mlop{"get"}, \mlval{"class"}] == \mlval{"motorway"}}\\\texttt{\mlprop{line-color}: \mlval{\#e892a2}}\\\texttt{\mlprop{line-width}: \mlval{3}}};
        \node[layer, below=0.5cm of l0] (l1) {\mlval{line}\\\texttt{\mlprop{filter}: [\mlop{"get"}, \mlval{"class"}] == \mlval{"primary"}}\\\texttt{\mlprop{line-color}: \mlval{\#fcd6a4}}\\\texttt{\mlprop{line-width}: \mlval{2}}};
        \node[layer, below=0.5cm of l1] (l2) {\mlval{line}\\\texttt{\mlprop{filter}: [\mlop{"get"}, \mlval{"class"}] == \mlval{"secondary"}}\\\texttt{\mlprop{line-color}: \mlval{\#f7fabf}}\\\texttt{\mlprop{line-width}: \mlval{1.5}}};

        \node[above=0.15cm of l0, font=\small\bfseries] {Input: 3 adjacent line layers};

        % Phase annotations (middle)
        \node[right=1.4cm of l0, check, text width=4.2cm] (p1) {%
            \stepcirc{1}~\textbf{Candidate check}\\[2pt]
            same type and source,
            source-layer differing props are feature-driven,
            no zoom ramps in diffs,
            no existing sort-key};
        \node[right=1.4cm of l1, check, text width=4.2cm] (p2) {%
            \stepcirc{2}~\textbf{Match detection}\\[2pt]
            all filters access the same property \texttt{\mlval{"class"}} with unique literals\\
            $\Rightarrow$ use \mlop{match} pattern};
        \node[right=1.4cm of l2, check, text width=4.2cm] (p3) {%
            \stepcirc{3}~\textbf{Zoom check}\\[2pt]
            all zoom ranges identical\\
            $\Rightarrow$ no guards needed};

        % Output (right)
        \node[merged, right=1.4cm of p2, text width=6.0cm] (out) {%
            \textbf{Merged \mlval{line} layer}\\[3pt]
            \texttt{\mlprop{filter}:}\\
            \texttt{~~[\mlop{"match"},\,[\mlop{"get"},\,\mlval{"class"}],}\\
            \texttt{~~~[\mlval{"motorway"},\mlval{"primary"},\mlval{"secondary"}],}\\
            \texttt{~~~\mlval{true},\,\mlval{false}]}\\[3pt]
            \texttt{\mlprop{line-color}:}\\
            \texttt{~~[\mlop{"match"},\,[\mlop{"get"},\,\mlval{"class"}],}\\
            \texttt{~~~\mlval{"motorway"},\,\mlval{"\#e892a2"},}\\
            \texttt{~~~\mlval{"primary"},\,\mlval{"\#fcd6a4"},\,\mlval{"\#f7fabf"}]}\\[3pt]
            \texttt{\mlprop{line-width}:}\\
            \texttt{~~[\mlop{"match"},\,[\mlop{"get"},\,\mlval{"class"}],}\\
            \texttt{~~~\mlval{"motorway"},\,\mlval{3},\,\mlval{"primary"},\,\mlval{2},\,\mlval{1.5}]}\\[3pt]
            \texttt{\mlprop{line-sort-key}:}\\
            \texttt{~~[\mlop{"match"},\,[\mlop{"get"},\,\mlval{"class"}],}\\
            \texttt{~~~\mlval{"motorway"},\,\mlval{0},\,\mlval{"primary"},\,\mlval{1},\,\mlval{2}]}
        };

        \node[above=0.15cm of out, font=\small\bfseries] {Output: 1 merged layer};

        % Arrows
        \draw[arrow] (l0.east) -- (p1.west);
        \draw[arrow] (l1.east) -- (p2.west);
        \draw[arrow] (l2.east) -- (p3.west);
        \draw[arrow] (p1.east) -- (out.west |- p1.east);
        \draw[arrow] (p2.east) -- (out.west);
        \draw[arrow] (p3.east) -- (out.west |- p3.east);
    \end{tikzpicture}}
    \caption{
Layer merging example.  Three adjacent line layers differing only in filter literal, color, and width are merged via the \mlop{match} pattern (Phase~2).
    The merged layer uses \texttt{[\mlop{"match"}, [\mlop{"get"}, \mlval{"class"}], \ldots]} for the filter, each differing paint property, and a synthezised sort key that preserves the original draw order.
    Phase~1 validates that all differing properties support feature-driven expressions, and Phase~3 confirms identical zoom ranges, so no zoom guards are needed}
    \label{fig:layer_merge_example}
\end{figure}

\subsection{Data-Level Optimizations}\label{sub:data_optimizations}

Data-level optimizations transform the tile data itself, reducing transfer size and decoding cost without altering the rendered output (constraint to \cref{eq:visual_equivalence}).
The encoding techniques used are surveyed in \cref{sec:data_encoding}.
\cref{tab:data_passes} summarizes the techniques we implemented.

\begin{table}[htbp]
\centering
\begin{tblr}{
  colspec = {p{3.9cm} p{1.8cm} p{9.5cm}},
  rows = {font=\small},
  row{1} = {font=\small\bfseries},
}
\toprule
Pass & Data & Description [compiler analogy] \\
\midrule
Tile shaving             & full scan & Encode only the properties, layers, and geometry types that the style actually references [tree shaking] \\
Transparent re-encoding  & full scan & Re-encode \ac{MVT} tiles into the \ac{MLT} format [instruction set translation] \\
Compression optimizations & full scan & Apply column-aware compression strategies based on data characteristics [\ac{PGO} optimizations] \\
\bottomrule
\end{tblr}
\caption{
Data-level optimization techniques, all requiring a full tile scan.
Together they form a logical, physical, and cost pipeline: tile shaving eliminates unreferenced data, re-encoding transforms the wire format, and compression optimizations select the most compact per-column encoding strategy.}
\label{tab:data_passes}
\end{table}

The pipeline operates in three stages that mirror a database query optimizer's \emph{logical} rewriting, \emph{physical} strategy selection, and \emph{cost-based} competition.

\begin{enumerate}
    \item \textbf{Style-driven tile shaving} (\cref{sub:tile_shaving}): the pipeline analyses the optimized style to produce a \emph{tile pruning advisory} that specifies, per source-layer and zoom level, which properties, geometry types, and property values to retain.
    It then prunes tiles according to this advisory before encoding.
    \item \textbf{Data transformations} (\cref{sub:string_property_interning,sub:feature_reordering}): the pipeline transforms remaining data to improve compressibility, replacing categorical string properties by frequency-ordered integers (\emph{string interning}) and reordering features to exploit spatial or categorical locality.
    \item \textbf{Encoding strategy selection} (\cref{sub:encoding_strategies}): the pipeline re-encodes the pruned, transformed data into \ac{MLT} using a multi-strategy competition that automatically profiles each data stream and selects the most compact encoding.
\end{enumerate}

\noindent The three stages compose:
shaving narrows the column set fed to Stage~2, which in turn determines the value distributions Stage~3 sees.
Stage~3's profiler-driven candidate pruning responds to those distributions $\bar{S}$ automatically.
Its viability tests (rle is considered if $\bar{S}_\text{rle} \geq 2.0$, delta bit-widths beat raw $\bar{S}_\Delta \geq 2.0$) operate on actual stream statistics, so a column whose distinct-value count drops after shaving naturally presents longer runs and triggers more compact candidates.

\paragraph{Style-Driven Tile Shaving}\label{sub:tile_shaving}

The bridge between style-level and data-level optimizations is the \emph{tile pruning advisory}, a structured report that captures, for each source-layer and zoom level, exactly what the rendering style needs.
We compute the advisory from the \emph{optimized} style (after all expression and structural passes) together with tile statistics.
The ordering matters because dead layer elimination, metadata refinement, and layer merging shrink the live-layer set, which in turn reduces the referenced properties and geometry types the advisory carries into the data level.
The empirical character of this style-to-data cascade is taken up in \cref{sub:mlt_results}.
The cascade mirrors the materialized-view maintenance problem, with the rendering style as the view definition, the tile set as the base relation, and the shaved tile as the materialized view.

\noindent The advisory captures the following per source-layer:

\begin{description}
    \item[\textbf{Used properties with zoom ranges}] for each property name appearing in a \texttt{[\mlop{"get"},~\ldots]} or \texttt{[\mlop{"has"},~\ldots]} expression in any live style layer, the zoom range over which it is referenced.
    Properties not referenced at any zoom level are candidates for removal.
    \item[\textbf{Used geometry types with zoom ranges}] which of Point, LineString, and Polygon are needed, derived from the style layer types (e.g., \texttt{\mlval{fill}} $\to$ Polygon, \texttt{\mlval{circle}} $\to$ Point).
    The pass removes features with geometry types unused by any style layer at the current zoom level.
    \item[\textbf{Unused zoom levels}] zoom levels at which no style layer targets the source-layer.
    Tiles at these zoom levels need not be encoded at all.
    \item[\textbf{Unused property values}] specific values of categorical properties that are never matched by any style filter.
    For example, if no style layer allows for a line with a \texttt{class=ferry} property, the pass can remove features with that value.
    \item[\textbf{Feature ID requirement}] whether any style layer uses \texttt{[\mlop{"id"}]} or \texttt{[\mlop{"feature-state"},~\ldots]}.
    When IDs are unused, the pass strips the entire ID column.
\end{description}

\noindent The pipeline applies the pruning advisory to each \ac{MVT} tile before encoding (\cref{fig:tile_shaving_pipeline}).

\begin{figure}[htbp]
    \centering
    \begin{tikzpicture}[
        box/.style={rectangle, draw, rounded corners=2pt, minimum height=0.65cm, font=\small, align=center},
        style/.style={box, fill=roleStyle!20, minimum width=2.8cm},
        data/.style={box, fill=roleData!25, minimum width=2.8cm},
        both/.style={box, fill=roleBoth!20, minimum width=2.8cm},
        arrow/.style={-Stealth, thick},
        lbl/.style={font=\scriptsize\itshape, text=darkgray, align=center},
        steps/.style={rectangle, draw, dashed, rounded corners=2pt, font=\small,
                      align=left, inner sep=6pt, text=black!80},
    ]
        % Row 1: style-side and data-side inputs
        \node[style] (sopt) {Style optimizations\\{\scriptsize (\cref{sub:expression_optimizations,sub:structural_optimizations})}};

        % Row 2: advisory
        \node[both, right=3cm of sopt] (adv) {Tile pruning advisory};
        \draw[arrow] (sopt) -- node[above, lbl] {optimized style} (adv);

        % Row 3:
        \node[data, right=3cm of adv] (stats) {Tile statistics};
        \draw[arrow] (stats) -- node[above, lbl] {property/\\geom distribution} (adv);

        % Row 3: MVT -> prune -> MLT
        \node[data, below=1.0cm of adv] (prune) {Prune \& transform};
        \node[data] at ($(sopt |- prune)$) (mvt) {\ac{MVT} tile};
        \node[data] at ($(stats |- prune)$) (mlt) {\ac{MLT} tile};

        \draw[arrow] (adv) -- (prune);
        \draw[arrow] (mvt) -- (prune);
        \draw[arrow] (prune) -- node[above, lbl] {re-encode} (mlt);

        % Pruning steps annotation
        \node[steps, below=0.5cm of prune] (ops) {%
            \textbf{Pruning operations (in order):}\\[2pt]
            1.\;Remove entire source-layers not referenced by any style layer.\\
            2.\;For each remaining layer:\\
            \quad a.\;Filter features by geometry type at current zoom.\\
            \quad b.\;Filter features by unused categorical property values.\\
            \quad c.\;Strip unused properties from the key-value table.\\
            \quad d.\;Strip feature IDs when not needed.%
        };
        \draw[dashed, gray] (prune) -- (ops);
    \end{tikzpicture}
    \caption{
Style-driven tile shaving pipeline.  The optimized style and tile statistics jointly produce a pruning advisory.  Each \ac{MVT} tile is pruned according to the advisory, transformed, and re-encoded as \ac{MLT}.  Style-level dead layer elimination narrows the advisory, which can reduce the data that must be encoded}
    \label{fig:tile_shaving_pipeline}
\end{figure}

\noindent After pruning, two data transformations improve compressibility before encoding.

\paragraph{String property interning}\label{sub:string_property_interning}
We replace categorical string properties (such as road class, land-use type, or administrative level) by frequency-ordered integer indices.
We derive the interning table from the tile statistics: we sort all distinct values of a string property by descending frequency and assign each value a sequential integer starting from~0.
The most common value receives index~0, the second most common receives~1, and so on.

This transformation has two benefits.
First, it reduces the per-feature storage cost from a variable-length string (with length prefix and dictionary overhead) to a small integer (typically \qtyrange{1}{2}{\byte} as a \ac{VarInt}, see \cref{sub:physical_integer_compressions}).
Second, the frequency ordering means that common values receive small indices, which compress well under \ac{VarInt} and delta encoding.

The optimizer rewrites the corresponding style expressions to reference the integer indices rather than the original string values.
For example, the optimizer rewrites a filter \texttt{[\mlop{"=="}, [\mlop{"get"}, \mlval{"class"}], \mlval{"primary"}]} to \texttt{[\mlop{"=="}, [\mlop{"get"}, \mlval{"class"}], \mlval{0}]} if \texttt{\mlval{"primary"}} is the most frequent value.
This rewriting ensures that the optimized style remains consistent with the transformed tile data.
Because interned tiles encode property values as style-specific integer indices rather than original strings, they are only compatible with the paired optimized style.
Other consumers (debugging tools, alternative renderers) would see opaque integers.
We therefore treat interning as an optional pipeline stage that can be disabled for multi-consumer deployments where tiles must remain self-describing.

\paragraph{Feature reordering}\label{sub:feature_reordering}
The order of features within a tile layer affects the compressibility of every parallel column.
The encoder explores multiple feature orderings (spatial via Hilbert/Morton curve of first vertex, ID-based, and property-based via grouping by a low-cardinality attribute) and retains the ordering that produces the smallest encoded output.
We describe this in detail in \cref{sub:encoding_strategies}.

\subsection{Reencoding data}

The encoder re-encodes the pruned and transformed data into the \ac{MLT} columnar format through a multi-strategy competition over per-column encodings, complemented by dictionary sharing across columns, geometry- and feature-ID-specific encoders, and a round-trip correctness check against the input.

\paragraph{Automatic Encoding Strategy Selection}\label{sub:encoding_strategies}
The re-encoding proceeds in three stages.
First, the encoder converts the row-oriented \ac{MVT} features into an owned columnar representation, with separate arrays for each property column, the geometry vertex buffer, and the feature ID column.
Second, it optionally reorders the columnar data according to the sort strategy competition.
Third, it independently encodes each column using the per-stream encoding competition.
The final wire format consists of three concatenated sections per layer: a header (layer name, tile extent, column count), a metadata section (column types, encoding descriptors), and a data section (the encoded streams).

We reformulate the encoding problem in our encoder as a \emph{multi-strategy competition} inspired by \ac{PGO} in compilers~\cite{changUsingProfileInformation1991} and the instance-optimized systems vision of SageDB~\cite{dingSageDBInstanceOptimizedData2022}:
rather than committing to a single encoding plan, the encoder uses the data distribution of each tile to try multiple strategies, measures the output size of each, and selects the smallest result.

This competition operates at two nested levels.
An \emph{outer loop} that tries multiple feature sort orderings, and an \emph{inner loop} that selects the best integer encoding for each individual data stream.
We deliberately make the inner competition criterion \emph{single-objective} (smallest encoded size).
We do not resolve multi-axis trade-offs (between transfer size, decoding latency, rendering throughput, and preprocessing cost) inside the encoder.
We report them separately per metric in the ablation evaluation (\cref{ch:experiments}), consistent with \cref{sec:problem_statement}: the deployment metrics of \cref{sub:metrics} are characterised per-axis rather than aggregated into a single objective, and the encoder's local size minimisation is a contributing quantity to the load-time axis rather than a global cost function.

A no-copy alternatives mechanism lets each candidate write into the encoder's buffers and rolls losing candidates back by resetting cursor positions, so the inner and outer competitions nest without redundant allocation.
Combined with warm-start caching of \ac{FSST} compressors and vertex-strategy decisions across sort trials, this keeps the cost of the nested competition near-additive rather than multiplicative in the number of trials. \cref{ch:experiments} reports the resulting preprocessing time empirically.
\Cref{fig:arch_enc} summarizes the encoder and decoder data-flow in \texttt{mlt-core}.

\begin{figure}[htbp]
    \centering
    \begin{subfigure}[t]{0.45\linewidth}
        \centering
        \begin{tikzpicture}[
            box/.style={rectangle, draw, rounded corners=2pt, minimum height=0.7cm,
                        minimum width=2.4cm, font=\small, align=center},
            decode/.style={box, fill=roleSrc!20},
            dec/.style={box, fill=roleOpt!35, minimum width=1.6cm},
            io/.style={font=\small},
            arrow/.style={-Stealth, thick},
            lbl/.style={font=\scriptsize\itshape, text=darkgray},
            node distance=1.7cm,
        ]
            % Main pipeline on centre line (x=0)
            \node[io]     (u8)     {Wire bytes};
            \node[decode, below=of u8, minimum width=3.2cm, align=center]     (name)
                {Columnar layer\\{\scriptsize per column: raw bytes}\\{\scriptsize or parsed values}};
            \draw[arrow] (u8)   -- node[left, lbl, align=right] {zero-copy parse} (name);
            \node[decode, below=of name] (tile)   {Row-oriented features};
            \draw[arrow] (name) -- node[left, lbl, align=right] {convert to\\row format} (tile);

            % Decoder off-centre towards the middle (right side)
            \node[dec, right=3.5cm of name.center, anchor=center] (decnode)
            {Decoder\\{\scriptsize reusable buffers,}\\{\scriptsize memory budget}\\{\scriptsize lazy decoding}};
            \draw[-Stealth, thick, densely dashed]
            ($(name.north) + (1,0)$) -- ++(0,0.8) -| (decnode.north);
            \draw[-Stealth, thick, densely dashed]
            (decnode.south) -- ++(0,-0.8) -| ($(name.south) + (1,0)$);

            % Phantom to balance bounding box
            \path (-4, 0);

            % Decoder arrows
        \end{tikzpicture}
        \caption{
Decoding pipeline.
An input byte slice is zero-copy parsed into a columnar layer whose columns are held lazily as either raw bytes or parsed values, decoded on demand.
The columnar data is converted to row-oriented features as needed.}
        \label{fig:arch_dec_sub}
    \end{subfigure}
    \hfill
    \begin{subfigure}[t]{0.45\linewidth}
        \centering
        \begin{tikzpicture}[
            box/.style={rectangle, draw, rounded corners=2pt, minimum height=0.7cm,
                        minimum width=2.0cm, font=\small, align=center},
            encode/.style={box, fill=roleData!25},
            enc/.style={box, fill=roleOpt!35, minimum width=1.6cm},
            io/.style={font=\small},
            arrow/.style={-Stealth, thick},
            lbl/.style={font=\scriptsize\itshape, text=darkgray},
            node distance=1.2cm,
        ]
            % Main pipeline on centre line (x=0)
            \node[io]     (write)   {Byte sink};
            \node[encode, below=of write]   (encoded) {Wire-ready bytes};
            \node[encode, below=1.4cm of encoded] (staged)  {Staged columnar layer};
            \node[encode, below=of staged]  (tile)    {Row-oriented features};

            % Encoder off-centre towards the middle (left side)
            \node[enc] at ($(staged)!0.5!(encoded) + (-2.5,0)$) (encnode) {Encoder\\{\scriptsize reusable buffers,}\\{\scriptsize caching}};

            \draw[arrow] (tile)    -- node[right, lbl, align=left] {convert to\\columnar} (staged);
            \draw[arrow] (staged)  -- node[right, lbl, align=left] {fully encode to\\wire-ready format} (encoded);
            \draw[arrow] (encoded) -- node[right, lbl] {zero-copy} (write);

            % Phantom to balance bounding box
            \path (3.5, 0);

            % Encoder arrows
            \draw[-Stealth, thick, densely dashed] (staged.west) -| (encnode.south);
            \draw[-Stealth, thick, densely dashed] (encnode.north) |- (encoded.west);
        \end{tikzpicture}
        \caption{
Encoding pipeline.
Row-oriented features are converted to a staged columnar layer, fully encoded to wire-ready byte buffers by the encoder that implements the multi-strategy competition, and written to a byte sink via a final zero-copy step.}
        \label{fig:arch_enc_sub}
    \end{subfigure}
    \caption{
Encoder~\subref{fig:arch_enc_sub} and decoder~\subref{fig:arch_dec_sub} data-flow in \texttt{mlt-core}}
    \label{fig:arch_enc}
\end{figure}

In an outer loop, the encoder runs a sort-strategy competition.
Reordering features within a layer permutes every parallel column (geometry, IDs, and all properties) simultaneously.
Different orderings favour different encoding schemes.
Spatial orderings cluster nearby geometries and their correlated properties, improving \ac{RLE} run lengths.
ID ordering produces sequential integer sequences amenable to delta-\acs{RLE}.
Which ordering wins depends on the value distributions and cardinalities of the specific layer.
A layer dominated by a single repeated property value may benefit most from spatial clustering, while a layer with sequential IDs and high-cardinality properties may favour ID ordering.
Because this balance shifts across layers, zoom levels, and tilesets, no single ordering dominates universally, and the choice cannot be resolved statically.

The encoder therefore tries multiple sort strategies on the same data and retains the encoding with the smallest total byte count; \cref{fig:sort_competition} diagrams the resulting trial-and-retain loop, including the bounding-box pruning that skips spatial trials for large, widely-spread layers.
The implemented strategies and their rationale are:

\begin{description}
    \item[\textbf{Unsorted}] the encoder preserves the original feature order.
    \item[\textbf{Spatial (Morton/Z-order)}] the encoder sorts features by the Z-order curve index of their first vertex, computed via bit-interleaving.  Fast to compute and provides reasonable spatial locality (\cref{sub:space_filling_curves}).
    \item[\textbf{Spatial (Hilbert)}] the encoder sorts features by the Hilbert curve index of their first vertex.  Slower to compute than Morton but achieves superior spatial locality, as the Hilbert curve avoids the long jumps inherent in the Z-order curve's recursive structure~\cite{moonAnalysisClusteringProperties2001}.
    \item[\textbf{ID}] the encoder sorts features by feature ID in ascending order.
\end{description}

For each candidate strategy, the encoder constructs a fresh \texttt{Encoder} instance, encodes the reordered layer, and measures the total output size (header + metadata + data).
The encoder retains the strategy producing the smallest output and discards all others.
When only a single strategy is applicable, the encoder elides the competition and encodes the layer directly without cloning.
We considered property-based sort strategies but did not include them in the current implementation, because the string interning pass (\cref{sub:string_property_interning}) already groups features of the same discriminator value into dense dictionary-index ranges, and Hilbert or ID ordering strictly dominated an explicit categorical sort on the benchmark tileset.

\begin{figure}[H]
    \centering
    \begin{tikzpicture}[
        node distance=0.5cm,
        box/.style={rectangle, draw, rounded corners=2pt, minimum width=4.2cm,
            minimum height=0.65cm, font=\small, align=center},
        input/.style={box, fill=roleInput!40},
        prep/.style={box, fill=roleMir!30},
        stage/.style={box, fill=roleStyle!20},
        encoder/.style={box, fill=roleData!25},
        decision/.style={diamond, draw, fill=roleDecision!35, aspect=2.8,
            font=\small, inner sep=1pt},
        output/.style={box, fill=roleOutput!25},
        arrow/.style={-Stealth, thick},
        label/.style={font=\scriptsize\itshape, text=darkgray},
        ]
        % Input
        \node[input] (layer) {\texttt{TileLayer01}};

        % Prep: group_string_properties
        \node[prep, below=of layer] (grp)
        {\texttt{group\_string\_properties()}};
        \draw[arrow] (layer) -- (grp);

        % Prep: collect_sort_strategies
        \node[prep, below=of grp] (coll)
        {\texttt{collect\_sort\_strategies()}};
        \draw[arrow] (grp) -- (coll);

        % Loop guard
        \node[decision, below=1.0cm of coll] (more) {more strategies?};
        \draw[arrow] (coll) -- (more);

        % Sort + columnarize
        \node[stage, below=1.0cm of more] (sort)
        {Sort + columnarize\\%
            {\scriptsize\texttt{from\_tile(clone, sort, groups)}}};
        \draw[arrow] (more) -- node[right, label] {yes} (sort);

        % Encode
        \node[encoder, below=of sort] (enc)
        {Encode all columns\\{\scriptsize (per-stream competition)}};
        \draw[arrow] (sort) -- (enc);

        % Compare
        \node[decision, below=of enc] (cmp) {size $<$ best?};
        \draw[arrow] (enc) -- node[right, label] {\texttt{total\_len()}} (cmp);

        % Yes: update best, then loop back
        \node[stage, below=0.5cm of cmp] (upd) {Update best};
        \draw[arrow] (cmp) -- node[right, label] {yes} (upd);

        % No: discard, loop back
        \draw[arrow] (upd.west) -| ($(more.west)+(-4cm,0)$) |- node[above, label] {no} (more.west);
        \draw[arrow] (cmp.west) -| ($(more.west)+(-4cm,0)$) |- (more.west);

        % Exit loop
        \node[output, right=2.2cm of more] (best) {Emit smallest encoding\\{\scriptsize\texttt{best.into\_layer\_bytes()}}};
        \draw[arrow] (more.east) -- node[above, label] {no} (best.west);
    \end{tikzpicture}
    \caption{
Multi-strategy sort competition.
        Shared-dictionary string groups are computed once and reused across all trials.
        The encoder collects candidate sort strategies, pruning spatial strategies via the bounding-box heuristic for large, widely-spread layers.
        For each candidate, the layer is cloned, sorted, and columnarized via \texttt{from\_tile}.
The columnar data is then encoded with per-stream integer competitions.
        The encoding with the smallest \texttt{total\_len()} is retained}
    \label{fig:sort_competition}
\end{figure}

To avoid wasting trials on unpromising strategies, a \emph{bounding-box heuristic} skips spatial sorting on large layers whose features are too spatially dispersed for locality clustering to help.
Specifically, for layers with more than \texttt{512} features, the heuristic prunes spatial trials when the first-vertex bounding box covers more than \texttt{80\%} of the tile extent on \emph{both} axes.
The encoder always tries small layers (below the threshold), since the absolute cost of an extra encoding pass is negligible for them.
We derived the thresholds experimentally from Germany-wide OpenMapTiles.

In an inner loop, a per-stream encoding competition tries every viable integer encoding for the data stream.
Each candidate pairs a \emph{logical} encoding (Plain, Delta, \ac{RLE}, or DeltaRle) with a \emph{physical} encoding (\ac{VarInt} or \ac{FastPFOR}~\cite{lemireDecodingBillionsIntegers2015}).
The general stream competition currently covers these four logical encodings.
The encoder applies geometry-specific strategies (Morton-delta, componentwise-delta) via a separate encoding path optimized for coordinate data, since their semantics are tied to the two-dimensional structure of vertex sequences and do not generalize to property columns.
Extending the general competition to include additional strategies (e.g., frame-of-reference without delta, or hybrid schemes) is a direction for future work, though the current candidate set captures the dominant encoding patterns observed in the benchmark tileset.
The naive approach of encoding every combination and comparing would be prohibitively expensive for large tiles.
Instead, a \emph{data profiler} samples a contiguous block taken from the middle of the stream: it profiles streams with at most $512$ values in full, and profiles longer streams over $\lfloor |s| / 100 \rfloor$ values clamped to $[512,\, 16384]$.
We chose the contiguous mid-stream window over a random sample because permutation destroys \ac{RLE} run lengths and delta-bit-width statistics, and because the middle of a spatially or ID-sorted stream is a representative point away from the boundary artefacts at the ends.
The profiler completes a single pass over the sample to prune candidates before any full encoding is attempted.
It excludes \ac{RLE} when the average run length $\bar{S} = |s| / n_{\text{runs}}$ falls below 2.0.
It excludes Delta when the stream is non-monotonic and the summed bit width of the zigzag-encoded deltas exceeds that of the raw values.
It excludes DeltaRle when the delta-stream average run length $\bar{S}_\Delta = (|s|-1) / n_{\Delta\text{-runs}}$ falls below 2.0 and neither \ac{RLE} nor Delta is independently viable.
The encoder only offers \ac{FastPFOR} for 32-bit integer types, as the algorithm's block structure is designed for fixed-width words.
The encoder fully encodes the remaining candidates (typically two to four) and commits the smallest output.

\paragraph{String encoding}
String columns compete across four strategies, reflecting the two orthogonal dimensions of string compression: deduplication and byte-level compression.

\begin{description}
    \item[\textbf{Plain}] raw bytes with length prefixes.
    The length stream is itself subject to the integer encoding competition.
    \item[\textbf{Dictionary}] a deduplicated corpus with per-string offset indices.
    Effective when many features share the same string value (e.g., road names that repeat across tiles).
    \item[\textbf{\ac{FSST}} (\cref{sub:string_compression})]~\cite{bonczFSSTFastRandom2020}
    A lightweight, random-access string compressor that learns a symbol table of up to~255 frequent byte sequences and replaces occurrences with single-byte codes.
    The encoder stores the symbol table once per column.
The compressed corpus remains directly scannable without decompression.
    \item[\textbf{\ac{FSST} + Dictionary}] \ac{FSST} applied to a deduplicated dictionary, combining both deduplication and byte-level compression.
\end{description}

A viability probe ensures that the encoder only attempts \ac{FSST} when the symbol-table overhead (up to \qty{2040}{\byte}) is likely to be recouped: it excludes columns below \num{2048}~raw bytes and evaluates a trial compression on a sample of up to \num{256} strings before committing.
The encoder caches trained \ac{FSST} compressors by column name and reuses them across all sort-strategy trials, avoiding redundant training.
The encoder applies a novel \emph{pruning variant} of \ac{FSST} that drops symbols from the persisted table when inlining their occurrences as escaped literals costs fewer bytes than retaining the table entry, shrinking the symbol table without changing the decoder.

\paragraph{Shared dictionary grouping}\label{sub:cardinality}
The \ac{MLT} format supports shared dictionaries where multiple string columns reference a common deduplicated corpus, amortising dictionary overhead across columns with overlapping vocabularies.
The reference encoder groups columns by naming prefix (e.g., \texttt{\mlval{name}}, \texttt{\mlval{name:en}}, \texttt{\mlval{name:de}}), which misses semantic overlaps (e.g., \texttt{\mlval{class}} and \texttt{\mlval{subclass}} in OpenMapTiles share many values).

\begin{sloppypar}
    Our encoder replaces prefix matching with MinHash-based similarity clustering (\cref{sub:lsh})~\cite{broderResemblanceContainmentDocuments1998}.
    \Cref{fig:minhash_clustering} walks through the four-step pipeline from per-column vocabularies through MinHash signatures and pairwise Jaccard estimation to the final Union-Find clustering.
    For each string column, the encoder computes two MinHash sketches of $128$ permutations: one over exact string values and one over byte trigrams (3-byte sliding windows).
    The choice of 128 permutations corresponds to an expected standard error of approximately $1/\sqrt{128} \approx 8.8\%$ on the Jaccard estimate, which is adequate for the binary accept/reject decision of the clustering step.
    The trigram sketch captures partial overlap between columns whose values are substrings or morphological variants of each other.
    The encoder estimates pairwise Jaccard similarity as $\hat{J} = |\{i : h_i^{(a)} = h_i^{(b)}\}| / 128$ for each sketch type and takes the maximum of two estimates.
    The encoder unifies columns whose similarity exceeds \texttt{7.5\%} via a Union-Find structure (union by size).
The low threshold reflects that even small vocabulary overlap reduces the marginal cost of adding a column to an existing dictionary.
    The encoder discards groups with fewer than two members.
    We selected the threshold and similarity mode via a sweep over $[2.5\%, 20\%]$ on the Germany OpenMapTiles tileset under three similarity modes (\emph{hybrid} taking $\max(\hat{J}_\text{exact}, \hat{J}_\text{trigram})$, \emph{only-trigram}, \emph{only-exact}), summarised in \cref{fig:minhash_sweep} in the appendix.
    Hybrid mode wins at every threshold at no significant performance cost, and the curve is nearly flat below \qty{7.5}{\percent}, where the empirical optimum at \qty{2.5}{\percent} saves only an additional \qty{1.6}{\mega\byte} on the \qty{3.14}{\giga\byte} tileset at the cost of false-positive union risk on tilesets with less correlated vocabularies.
    We therefore retain \qty{7.5}{\percent} as a conservative default that captures the majority of the grouping benefit.
\end{sloppypar}

\DeclareRobustCommand{\circled}[1]{\tikz[baseline=(char.base)]{\node[circle, draw, inner sep=1pt, font=\scriptsize\bfseries] (char) {#1};}}

\begin{figure}[htbp]
    \centering
    \resizebox{\textwidth}{!}{\begin{tikzpicture}[
            arrow/.style={-Stealth, thick},
            cmp/.style={font=\scriptsize, align=left},
            stepbox/.style={draw, rounded corners=3pt, inner sep=5pt, font=\scriptsize},
            steplabel/.style={font=\scriptsize\bfseries},
            ]
            % Top-left (Step 1): Column value table
            \node[stepbox] (tbl) {%
                \begin{tabular}{@{}l l@{}}
                    \textbf{Column} & \textbf{Distinct values (sample)} \\
                    \midrule
                    \texttt{\mlval{name:de}}  & Straßburg, Freiburg, Mühlhausen, \ldots \\
                    \texttt{\mlval{name:fr}}  & Strasbourg, Fribourg, Muhlhouse, \ldots \\
                    \texttt{\mlval{class}}    & motorway, trunk, primary, \ldots \\
                    \texttt{\mlval{subclass}} & motorway, trunk, primark, \ldots \\
                    \texttt{\mlval{maxspeed}} & 30, 50, 80, \ldots
                \end{tabular}%
            };
            \node[steplabel, above=1pt of tbl] {\circled{1} Column vocabularies};

            % Top-right (Step 2): MinHash sketch generation
            \node[stepbox, right=0.8cm of tbl, anchor=west] (hash) {%
                \begin{tabular}{@{}l c c@{}}
                    \textbf{Column} & \textbf{Exact sketch} & \textbf{Trigram sketch} \\
                    \midrule
                    \texttt{\mlval{name:de}}  & $[h_1\!\ldots\!h_{128}]_{\text{val}}$ & $[h_1\!\ldots\!h_{128}]_{\text{3g}}$ \\
                    \texttt{\mlval{name:fr}}  & $[h_1\!\ldots\!h_{128}]_{\text{val}}$ & $[h_1\!\ldots\!h_{128}]_{\text{3g}}$ \\
                    \texttt{\mlval{class}}    & $[h_1\!\ldots\!h_{128}]_{\text{val}}$ & $[h_1\!\ldots\!h_{128}]_{\text{3g}}$ \\
                    \texttt{\mlval{subclass}} & $[h_1\!\ldots\!h_{128}]_{\text{val}}$ & $[h_1\!\ldots\!h_{128}]_{\text{3g}}$ \\
                    \texttt{\mlval{maxspeed}} & $[h_1\!\ldots\!h_{128}]_{\text{val}}$ & $[h_1\!\ldots\!h_{128}]_{\text{3g}}$ \\
                \end{tabular}%
            };
            \node[steplabel, above=1pt of hash] {\circled{2} MinHash signatures (128 permutations)};

            % Bottom-right (Step 3): Pairwise Jaccard
            \node[stepbox, anchor=north] (sim) at ($(hash.south)+(0,-0.8cm)$) {%
            	\begin{tabular}{@{}l c c l@{}}
            		\textbf{Column pair} &
            		$\hat{J}_{\text{exact}}$ &
            		$\hat{J}_{\text{tri}}$ &
            		\textbf{Observation} \\
            		\midrule
            		\texttt{\mlval{name:de}} vs.\ \texttt{\mlval{name:fr}}
            		& 0.14 & \textbf{0.35}
            		& shared trigrams\\

            		\texttt{\mlval{class}} vs.\ \texttt{\mlval{subclass}}
            		& 0.45 & \textbf{0.57}
            		& shared values \\

            		\texttt{\mlval{name:de}} vs.\ \texttt{\mlval{class}}
            		& 0.00 & \textbf{0.01}
            		& unrelated \\

            		\texttt{\mlval{maxspeed}} vs.\ others
            		& 0.00 & 0.00
            		& no overlap \\
            	\end{tabular}%
            };
            \node[steplabel, below=1pt of sim] {\circled{3} Pairwise comparison ($\hat{J} = |\{i : h_i^{(a)} = h_i^{(b)}\}| \,/\, 128$)};

            % Bottom-left (Step 4): Union-Find result
            \node[stepbox, cmp] at ($(tbl |- sim)$)  (res) {%
                \textbf{Shared dictionary groups:}\\[3pt]
                \fbox{\strut\texttt{\mlval{name:de}}, \texttt{\mlval{name:fr}}}\\[1pt]
                \quad{\scriptsize trigrams detected cross-language overlap}\\[4pt]
                \fbox{\strut\texttt{\mlval{class}}, \texttt{\mlval{subclass}}}\\[1pt]
                \quad{\scriptsize exact values detected semantic overlap}\\[4pt]
                \texttt{\mlval{maxspeed}} - individual dictionary\\[1pt]
                \quad{\scriptsize no overlap with any string column}%
            };
            \node[steplabel, below=1pt of res] {\circled{4} Union-Find clustering};

            % Arrows (clockwise)
            \draw[arrow] (tbl) -- node[above, font=\scriptsize\itshape] {hash} (hash);
            \draw[arrow] (hash) -- (sim);
            \draw[arrow] (sim) -- node[above, font=\scriptsize\itshape] {$>\tau$} (res);

    \end{tikzpicture}}
    \caption{
MinHash-based shared dictionary clustering (clockwise from top-left).
        \circled{1}~Each string column's distinct values form a set.
        \circled{2}~Two 128-permutation MinHash sketches are computed per column: one over exact string values, one over byte trigrams (3-byte sliding windows).
        \circled{3}~Pairwise Jaccard similarity is estimated by counting matching hash positions.
The maximum of the two sketch types is compared against threshold $\tau = 0.075$.
        Exact matching detects columns that share literal values (\texttt{\mlval{class}}/ \texttt{\mlval{subclass}}).
Trigram matching detects cross-language morphological variants (\texttt{\mlval{name:de}}/ \texttt{\mlval{name:fr}} in a border region).
        \circled{4}~Column pairs exceeding the threshold are unified via Union-Find.
The resulting groups share a single deduplicated dictionary corpus}
    \label{fig:minhash_clustering}
\end{figure}

\paragraph{Geometry encoding}
The encoder encodes geometry vertices as either \emph{componentwise deltas} (CWD), \emph{Hilbert-code dictionaries}, \emph{Morton-code dictionaries} or plain vertices.
CWD encodes each vertex as the delta from the previous vertex, producing a single interleaved stream $[x_0, y_0, \Delta x_1, \Delta y_1, \ldots]$ that is then compressed via the per-stream integer competition.
The Morton/Hilbert-code dictionary approach computes a Z-order code for each unique vertex pair, stores the deduplicated codes as a sorted dictionary, and references them via per-vertex offset indices, producing two streams (offsets and dictionary deltas) instead of one.
The encoder tries both strategies (when the coordinate range permits Morton encoding, i.e., both axes fit within 16~bits after shifting) and retains the smaller result.


\paragraph{Feature IDs encoding}
The encoder encodes feature IDs through the same per-stream integer competition as other integer columns, with an explicit tag selecting between 32-bit and 64-bit word widths.
The encoder narrows 64-bit-declared ID columns to 32 bits when every observed value fits.

The same data profiler that drives the inner competition discovers structural patterns.
Sequential IDs reach DeltaRle, constant IDs reach \ac{RLE}, and short or irregular sequences fall back to plain \ac{VarInt}.

\paragraph{Encoder correctness validation}\label{par:encoder_validation}
Because \ac{FastPFOR} disclaims input validation and the encoder's narrowing, run-encoding, and dictionary rebuild stages have independent failure modes, we validate correctness through property-based round-trips per integer encoder, coverage-guided fuzzing over full layer round-trips, snapshot tests over the profiler's candidate selection, and the pixel-exact end-to-end check from \cref{sec:visual_correctness}.
This soft-correctness guarantee underwrites the experimental results in \cref{ch:experiments}.

\subsection{Tile Statistics Collection}\label{sub:tile_statistics}

Several optimizations passes require knowledge of the actual tile data, namely stats-driven constant folding and selectivity reordering (\cref{sub:expression_optimizations}), dead elimination with geometry-type mismatch and metadata refinement (\cref{sub:structural_optimizations}), and the tile pruning advisory (\cref{sub:tile_shaving}).
A statistics collection phase provides this knowledge by scanning the tile set and producing a per-source, per-source-layer statistical profile.
The optimizer accepts either an MBTiles database, in which case the scan runs on demand, or a pre-computed statistics file produced by an earlier scan.

For each source layer, the phase collects the following statistics:

\begin{description}
    \item[\textbf{Feature count}] total number of features across all tiles\footnote{Contrairy to na"ive reasoning, eatures appearing in multiple tiles at different zoom levels are counted multipe times, as the optimizer reasons per-tile, not per-geographic-entity.
This creates bias, but bias that is representative to the usecase of rendering maps, where tiled features are used instead of geographic features.}.
    \item[\textbf{Features by zoom}] a per-zoom-level feature count, used by metadata refinement to determine the first and last zoom level at which data exists.
    \item[\textbf{Geometry types}] counts of Point, LineString, and Polygon features, used by dead elimination to detect geometry-type mismatches.
    \item[\textbf{Feature ID presence}] whether any feature carries a non-zero feature ID, used by the advisory to decide whether to strip the ID column.
    \item[\textbf{Per-property statistics}] for each property name, the statistics depend on the wire type:
    \begin{description}
        \item[\emph{Boolean}] present count and true count.
        \item[\emph{Integer/unsigned integer}] present count, minimum, maximum, cardinality, and (if cardinality $\leq 200$) a full value-frequency map.
        \item[\emph{Double}] present count, minimum, maximum, cardinality (no value-frequency map, since floating-point equality is unreliable for folding).
        \item[\emph{String}] present count, cardinality, and (if cardinality $\leq 200$) a value-frequency map sorted by descending frequency.
    \end{description}
\end{description}

We allow collecting statistics at a configurable \emph{sample rate}.
A sample rate of~1.0 (full census) is required for irreversible decisions such as dead layer elimination or property interning.
Lower sample rates are sufficient for selectivity estimation and reordering.
A sample-rate guard in the optimizer enforces this distinction.
Passes that permanently remove data check that the sample rate is~1.0 before applying.

\subsection{Preprocessing Cost Model}\label{sub:preprocessing_cost_model}

All optimizations described above run \emph{offline}, paid once per (style, tileset) build and amortised across every subsequent tile request and rendering session.
The total preprocessing cost decomposes into four additive terms covering style passes (\cref{sub:expression_optimizations,sub:structural_optimizations}), tile statistics (\cref{sub:tile_statistics}), advisory computation (\cref{sub:tile_shaving}), and \ac{MLT} re-encoding (\cref{sub:encoding_strategies}).
\begin{equation}\label{eq:cost_model}
    C_\text{total}(S, T) \;=\; C_\text{style}(S) \;+\; C_\text{stats}(T) \;+\; C_\text{adv}(S) \;+\; C_\text{enc}(S, T).
\end{equation}
Bounding-box pruning, data profiling, \ac{FSST} warm-start caching, scratch-buffer recycling, and the \ac{RAII} alternatives guard together keep $C_\text{enc}$ bounded under the multi-strategy competition.
We report wall-clock cost on the Germany OpenMapTiles tileset in \cref{sub:preprocessing_cost}.
The focus of this thesis is on the delivery and rendering improvements that this one-off cost buys.

\section{Gathering a Representative Sample}\label{sec:gathering_sample}

Our deep evaluation runs 15 styles (\cref{sub:deep_corpus}) against reproducible tile data (\cref{sub:tile_data}) and a stratified scenario sample (\cref{sub:scenarios}).
A further 13 styles (\cref{sub:generalization_corpus}) carry static-only checks to test whether gains generalise beyond the deep tier.

\subsection{Benchmarking Corpus}\label{sub:deep_corpus}

We curated a corpus of 15 publicly available styles that use the OpenMapTiles or BasemapDE schema:
four \textbf{OpenFreeMap} styles (liberty, bright, positron, fiord);
seven \textbf{OpenMapTiles community} styles (dark-matter, osm-bright, klokan-basic, toner, osm-liberty, americana, stadia-outdoors);
and four \textbf{government and institutional} styles (ICGC~fosc, ICGC~gris, Basemap\-DE\footnote{\url{https://basemap.de/}, last accessed 2026-05-23} topo\-graphic, Basemap\-DE color).

These styles span a range of complexity-from minimal basemaps with fewer than 50~layers to verbose institutional styles with several hundred layers-and exercise different subsets of the expression language.

\subsection{Generalisation Corpus}\label{sub:generalization_corpus}

To assess whether optimizations gains generalize beyond the deep corpus, we collected 13 styles from the Maputnik style editor catalog (two additional catalog entries were unavailable at the time of measurement).
These styles are evaluated using static analysis only (style size, complexity metrics, pass applicability) and do not require tile data.

\subsection{Tile Data and Caching}\label{sub:tile_data}

We source tile data from \ac{OSM} via Planetiler or the German government.
We consume both sources through the standard Web Mercator XYZ tiling scheme used by MapLibre, so all zoom-level reasoning in the thesis assumes this grid.
We access BasemapDE through its XYZ-compatible endpoint rather than its native \acs{WMTS}~\cite{masoOpenGISWebMap2010} grid, ensuring that the optimizer's tile indexing is uniform across the corpus.
Only for \ac{OSM} is the underlying data fully available, so only for this source can we collect tile statistics at a sample rate of~1.0.
To ensure reproducibility across benchmark runs, we route all tile requests through a local caching proxy that stores tiles on first access and replays them on subsequent runs with the latency and download speed of a 4G internet connection according to Rusan et al.~\cite{rusanAssessmentPacketLatency2015} ($\frac{\text{size}}{30\,\mathrm{MiB/s}} + 25ms$).
This eliminates network variability and guarantees that every ablation step operates on identical tile data (after an initial, discarded run).
The model does not simulate packet loss, HTTP/2 multiplexing, TCP slow-start, or HTTP/3 (QUIC) with its 0-RTT connection resumption and absence of head-of-line blocking.
All of these would amplify the benefit of smaller tiles, so the measured load-time improvements from tile shaving are a conservative lower bound on what a real deployment would observe.
This ensures that we measure tile-fetch improvements (e.g.\ smaller tiles due to shaving) under conditions representative of a \ac{CDN}-served deployment while remaining deterministic across runs.
The same network model also feeds the per-style \ac{CDN} break-even analysis in \cref{sub:mlt_discussion}, which extends this single-client setup to a multi-style cache.

\subsection{Geographic and Use-Case Scenarios}\label{sub:scenarios}

The benchmark harness defines 18 scenarios drawn from a matrix of 15 geographic locations and 5 camera animation types.
\textbf{locations} span high-density cities with Latin scripts (Munich, Berlin, Paris, New York, Rome), cities with non-Latin scripts (Tokyo, Beijing, Seoul, Cairo), and low-density or special nature regions (Black Forest, Amazon, Swiss Alps, Sahara, Venice, Stockholm Archipelago).
\textbf{animations} exercise different rendering workloads:
\emph{zigzag} alternates zoom levels while panning;
\emph{spiral} traces a circle at constant zoom with rotating bearing;
\emph{zoomdrill} zooms from overview to street level and back;
\emph{pansweep} sweeps across a region at constant zoom;
and \emph{bearingspin} rotates bearing at high pitch.
See \cref{app:animation_patterns} for the keyframe diagrams.

Each scenario is a specific (location, animation) pair.
The full cross-product of 15 locations and 5 animations is 75 pairs, which is too expensive\footnote{Each cell requires 15 repeat runs for statistical validity across 19 ablation steps and 15 styles, as this would result in $75 \times 15 \times 19 \times 15 = \num{320625}$ full animations} to run for every ablation step.
We chose the 18 selected scenarios by stratified sampling so that every location appears at least once and every animation appears at least three times, and so that all combinations of density class (high/low) and script class (Latin/non-Latin) are represented.
Within each stratum we selected the specific (location, animation) pair that exercises the most layers in the densest style in the corpus, ensuring that the sample is biased \emph{against} the null hypothesis of \enquote{no effect}, i.e.\ towards scenarios where any rendering-cost differences are most likely to be visible.

This is a convenience sample by design:
we do not claim it is unbiased across all real-world usage patterns, but we claim it is a reasonable stress test of the optimizer across the dimensions (density, script, zoom trajectory, camera motion) that most plausibly interact with style and data-level optimizations.

\section{Evaluation}\label{sec:evaluation}

Since many of the investigated optimizations techniques have not previously been studied in the context of map rendering pipelines, we must evaluate their effectiveness empirically.
We carry out this empirical evaluation in \cref{ch:experiments}.
The rest of this section defines the methodology that chapter applies.
Our evaluation reports the deployment metrics framed in \cref{sec:problem_statement} along several axes: data savings (which feed the load-time axis), rendering performance, optimization cost (the $C_\text{total}$ breakdown of \cref{sub:preprocessing_cost_model}), and qualitative efficiency.

\begin{description}
    \item[\textbf{Data savings}] reduction in style size (raw, gzip, and Brotli) and tile transfer volume.
    \item[\textbf{Rendering performance}] client-side metrics including load time, frames per second, frame time percentiles (p50, p95, p99), jank count, style parse time, time to first tile, and time to first frame.
    \item[\textbf{Optimization cost}] preprocessing time introduced by each pass, measured as wall-clock time on the optimizer's host machine.
This is an offline cost amortised over all subsequent tile deliveries.
    \item[\textbf{Efficiency}] \emph{qualitative} discussion of the energy implications of reduced data transfer and fewer rendering operations.  We perform no direct power or energy measurement (\ac{RAPL}, battery delta), because the benchmark harness runs on a desktop workstation rather than a thermaly and battery-constrained mobile device.
Any quantitative energy claim would require a different hardware platform and we defer it to future work.
\end{description}

\subsection{Metrics}\label{sub:metrics}

Per run, the benchmark harness collects the following metrics:

\begin{description}
    \item[\textbf{Load time}] wall-clock time from requesting the map to completion of the initial render, including style parsing, tile fetching, tile decoding, and first paint.
    Within the rendering benchmark, decoding latency is embedded in \texttt{loadMs} rather than isolated, because the browser-based harness cannot separate decoder wall-clock time from the surrounding fetch and first-paint work without renderer-internal instrumentation.
    We instead measure isolated decoding throughput via Criterion\footnote{\url{https://github.com/bheisler/criterion.rs}, last accessed 2026-05-23} micro-benchmarks on our decoder (\cref{tab:decode_throughput}).
    \item[\textbf{Style parse time}] time spent parsing and compiling the style \ac{JSON} into the renderer's internal representation.
    \item[\textbf{Time to first tile} (\texttt{firstTileMs})] time from map initialisation until the first tile is decoded and ready for rendering.
    \item[\textbf{Time to first frame}] time until the first frame is composited and presented.
    \item[\textbf{Frames per second} (\texttt{fps})] average rendering throughput during the animation phase.
    \item[\textbf{Frame time percentiles} (\texttt{p50}, \texttt{p95})] distribution of per-frame render times, capturing both typical and worst-case latency.
    \item[\textbf{Jank count}] number of frames exceeding \qty{16.67}{\milli\second} (the 60\,\ac{FPS} budget), indicating visible stuttering.
    \item[\textbf{Heap memory}] JavaScript heap usage during the benchmark, capturing the memory cost of style structures and tile data.
    \item[\textbf{Idle time}] time the renderer spends idle (waiting for tiles or \ac{GPU} sync), indicating headroom.
\end{description}

\noindent Additionally, we compute static metrics for each optimized style without running the renderer.
\textbf{Style size} covers raw byte count, gzip-compressed size, and Brotli-compressed size.
\textbf{Complexity metrics} comprise total \ac{AST} node count, maximum expression depth, layer count, filter count, and expression operator histogram.

\subsection{Ablation Methodology}\label{sub:ablation}

Our primary evaluation methodology is cumulative ablation: we add passes one at a time in a fixed order and re-run the full benchmark suite after each addition.
This reveals both the marginal contribution of each pass and the interactions between passes.
The ablation comprises 19~cumulative steps (a baseline with no optimizations followed by passes added incrementally in the order listed in \crefrange{tab:expression_passes}{tab:data_passes}) plus one additional configuration that applies plain tile shaving on top of the unoptimized style, included to isolate the style$\leftrightarrow$shaving interaction in \cref{sub:mlt_results}.
\cref{sub:combined_ablation,tab:marginal_contribution} present the step-by-step results.

We execute each scenario/style combination for 15~runs.
We report the median and standard deviation of each metric.
We prefer the median over the mean because rendering performance distributions are right-skewed (occasional long frames inflate the mean).
To reduce measurement noise, the benchmark harness pins rendering to real \ac{GPU} hardware via ANGLE\footnote{\url{https://github.com/google/angle}, last accessed 2026-05-23}/Vulkan\footnote{\url{https://www.vulkan.org/}, last accessed 2026-05-23} rather than the software rasteriser used for pixel-exact render tests.
This ensures that we measure \ac{GPU}-bound workloads (fragment shading, texture sampling, blend operations) realistically.

We serve all tile requests from a local caching proxy as described in \cref{sub:tile_data}.
The benchmark harness itself is a Puppeteer\footnote{\url{https://pptr.dev/}, last accessed 2026-05-23}-driven headless browser that programmatically loads the MapLibre GL~JS renderer, applies a camera animation scenario, and collects performance counters via the \texttt{Performance} \ac{API}\footnote{\url{https://developer.mozilla.org/en-US/docs/Web/API/Performance_API}, last accessed 2026-05-23} and custom instrumentation hooks in the renderer.

\begin{chaptersummary}
The optimizer is a \textbf{three-level pipeline} (\cref{sub:pipeline}) operating on the $(S, T) \to (S', T')$ problem of \cref{sec:problem_statement}: eight expression-level passes run to a fixpoint, six structural passes act on a typed style document, and data-level passes re-encode shaved tiles into \ac{MLT} via per-stream strategy competition.
The \textbf{pruning advisory} derived from the optimized style is what ties the levels together: a projection-pushdown plan over the tile data.
We measure the pipeline with a \textbf{19-step cumulative ablation} (plus one plain-tile-shaving comparison configuration) across \textbf{15~styles, 18~stratified scenarios}, and 15~repeat runs per cell on the Germany OpenMapTiles tileset, plus a 13-style generalization corpus and a preprocessing-cost model that decomposes the offline cost into four additive stages.

The following chapter applies this methodology to evaluate each optimization category in isolation, tests generalization beyond the benchmarking corpus, and analyses the full pipeline's combined effect.
\end{chaptersummary}
